\title{CARTA DE SERVIÇOS DO CRP~SP}\author{}\date{}

\begin{document}\maketitle

\section{Apresentação}

Com o intuito de informar a sociedade sobre os principais serviçosprestados pelo Conselho Regional de Psicologia de São Paulo,disponibilizamos a Carta de Serviços do CPR SP, material de amplo escopoe fácil acesso, construído tendo em vista compromisso com uma gestão dequalidade, voltada para o constante aprimoramento dos nossos serviços epara a garantia do exercício profissional da Psicologia de forma ética esocialmente comprometida.

Todos os serviços que o CRP~SP disponibiliza estão detalhados nestematerial, que inclui o objetivo precípuo da entidade e seus canais decomunicação e atendimento seja à categoria de psicólogas/os, seja àpopulação em geral.

\section{O CRP~SP (páginas 11 e 12 do Relatório deGestão)}

O Conselho Regional de Psicologia de São Paulo (6ª~Região) é umaentidade dotada de personalidade jurídica de direito público, comautonomia administrativa e financeira, nos termos da Lei Federalnº~5.766, de 20 de dezembro de 1971. Caracteriza-se como autarquia defiscalização do exercício e das atividades profissionais, constituindoserviço público federal, vinculado ao Conselho Federal de Psicologia(CFP) como parte do Sistema Conselhos, composto atualmente por 24conselhos regionais e o CFP.

O CRP~SP é o maior conselho regional do Sistema Conselhos, abarcandoquase 1/3 da categoria profissional de psicólogas/os, com cerca de147mil inscrições ativas de pessoas físicas e 12 mil de pessoas jurídicas.Organizado de maneira descentralizada em todo o estado de São Paulo, édividido em uma sede administrativa, na capital, e 11 subsedes:

\begin{itemize}
  \item  Alto Tietê (Mogi das Cruzes);
  \item  Assis;
  \item  Baixada Santista e Vale do Ribeira (Santos);
  \item  Bauru;
  \item  Campinas;
  \item  Grande ABC (Santo André);
  \item  Metropolitana (São Paulo);
  \item  Ribeirão Preto;
  \item  São José do Rio Preto;
  \item  Sorocaba;
  \item  Vale do Paraíba e Litoral Norte (Taubaté)
\end{itemize}

Além disso, também atuam no conselho 142 trabalhadoras/es efetivas/os,11 trabalhadoras/es em cargos de livre provimento, 27 estagiárias/os, 29prestadoras/es de serviços de limpeza e vigilância, 14 conselheiras/osefetivas/os, 15 conselheiras/os suplentes, 11 comissões gestoras(formadas por 52 gestoras/es) e mais de 380 psicólogas/os que colaboramem diversas comissões.

Missão

Atuar com eficácia na fiscalização do exercício da profissão depsicóloga/o, com a responsabilidade de orientar, disciplinar e zelarpela fiel observância dos princípios ético-profissionais, contribuindopara o desenvolvimento da Psicologia como ciência e profissão na buscado bem comum.

Visão

Ser reconhecido pela defesa de direitos humanos e sociais, promovendosaúde e qualidade de vida para as pessoas e coletividade comresponsabilidade social crítica.

Valores

Promover a liberdade, a dignidade, a igualdade e a integridade daspessoas e coletividades, contribuindo para a eliminação de quaisquerformas de negligência, crueldade e opressão, atuando para odesenvolvimento da Psicologia como campo científico de conhecimento eprática.

\section{Psicologia: como se estrutura e como seconecta}

\href{https://d.docs.live.net/c339283d8975e0ad/revisão/Fwd\%20Carta\%20de\%20Serviços\%20CRP\%20SP/Relatorio-de-gestao_2023_CRPSP-13.pdf}{Página13 do relatório de gestão}

\href{https://d.docs.live.net/c339283d8975e0ad/revisão/Fwd\%20Carta\%20de\%20Serviços\%20CRP\%20SP/Relatorio-de-gestao_2023_CRPSP-14.pdf}{Página14 do relatório de gestão}

\part{Serviços à pessoafísica}

\section{O que é?}

É o registro necessário para a atuação profissional em Psicologia. Oexercício profissional no estado de São Paulo estará autorizado após ageração do número de registro profissional.

\section{Quem deve se inscrever?}

Profissionais graduadas/os em Psicologia, que tenham atuação no estadode São Paulo.

\section{Como solicitar o serviço?}

Por meio do \emph{site} do CRP~SP, na guia\href{https://www.crpsp.org/pagina/view/308}{Atendimento/Pessoafísica}.

\section{Que informações ou documentos são necessários?}

\begin{itemize}
  \item  documento de identificação (RG, CNH etc.);
  \item  Cadastro de Pessoa Física (CPF);
  \item  certidão de quitação eleitoral;
  \item  comprovante ou declaração de endereço;
  \item  diploma de formação em Psicologia (frente e verso) ou certidão de  colação de grau;
  \item  certidão de casamento ou de união estável;
  \item  comprovante de quitação com o serviço militar (para homens até 45  anos)
\end{itemize}

\section{Quais são as etapas para a realização desse serviço?}

\begin{enumerate}
  \item  A/O psicóloga/o realiza o requerimento de inscrição \emph{on-line},  pelo \href{https://www.crpsp.org/pagina/view/308}{\emph{site} do  CRP~SP};
  \item  o CRP~SP faz a análise da documentação;
  \item  os boletos para pagamento da taxa de inscrição e da anuidade são  emitidos;
  \item  após o pagamento das taxas, o registro é deferido e uma declaração  para atuação profissional é disponibilizada.
\end{enumerate}

\section{Quanto custa?}

Os valores para 2024 ficaram estabelecidos assim:

\begin{itemize}
  \item  pagamento à vista: R\$~575,60;
  \item  pagamento parcelado: 6 x R\$~95,93
\end{itemize}

\section{Quanto tempo leva?}

Até 30 dias.

\section{Como ter acesso ao atual \emph{status} do serviço?}

Ao se fazer a solicitação é gerada uma senha de acesso. Com essa senha,é possível acompanhar o \emph{status} da solicitação \emph{on-line.}

\section{Legislação relacionada}

\begin{itemize}
  \item  \href{https://atosoficiais.com.br/cfp/resolucao-administrativa-financeira-n-3-2007-institui-a-consolidacao-das-resolucoes-do-conselho-federal-de-psicologia?origin=instituicao&q=resolu\%C3\%A7\%C3\%A3o\%203/2007}{Resolução  CFP nº~03/2007} (Consolidação das Resoluções do Conselho Federal de  Psicologia);
  \item  \href{https://atosoficiais.com.br/cfp/resolucao-administrativa-financeira-n-20-2018-altera-o-manual-de-procedimentos-administrativo-e-financeiro-do-sistema-conselhos-de-psicologia-anexo-da-resolucao-cfp-n-202018-a-resolucao-cfp-n-03-2007-a-resolucao-cfp-n-16-2019-e-da-outras-providencias}{Resolução  CFP nº~20/2018}, que determina a revisão e ampliação do  \href{https://transparencia.cfp.org.br/wp-content/uploads/sites/7/2020/05/Manual-de-Procedimentos-Administrativos-e-Financeiro-CFP-RES-20.2018.pdf}{\emph{Manual  de procedimentos administrativos e financeiros}}
\end{itemize}

\section{Serviços correlatos}

\begin{itemize}
  \item  \hyperref[]{Entrega  carteira de identidade profissional (CIP)};
  \item  \hyperref[]{alteração  de inscrição provisória para definitiva};
  \item  \hyperref[]{inscrição secundária};
  \item  \hyperref[]{transferência de  registro profissional};
  \item  \hyperref[]{cancelamento de  registro profissional};
  \item  \hyperref[]{prorrogação  do prazo de entrega do diploma};
  \item  \hyperref[]{solicitação  de isenção e interrupção temporária de anuidades}
\end{itemize}

\chapter{\texorpdfstring{\emph{Entrega da carteira de identidadeprofissional(CIP)}}{Entrega da carteira de identidade profissional (CIP)}}

\section{O que é?}

A CIP é o documento que identifica a/o profissional inscrita/o eregular. Nela consta o número de registro junto ao CRP~SP.

\section{Quem pode utilizar esse serviço?}

Profissionais graduadas/os em Psicologia, que tenham atuação no estadode São Paulo e que tenham realizado a inscrição no CRP~SP.

\section{Como solicitar o serviço?}

Realizar o \href{https://www.crpsp.org/agendamento}{agendamento préviono \emph{site} do CRP~SP} e apresentar, na data e local agendados, osoriginais dos documentos solicitados. Nessa ocasião serão coletados osdados biométricos.

\section{Que informações ou documentos são necessários?}

\begin{itemize}
  \item  documento de identificação (RG, CNH etc.);
  \item  Cadastro de Pessoa Física (CPF);
  \item  certidão de quitação eleitoral;
  \item  comprovante ou declaração de endereço;
  \item  diploma de formação em Psicologia (frente e verso) ou certidão de  colação de grau;
  \item  certidão de casamento ou de união estável;
  \item  comprovante de quitação com o serviço militar (para homens até 45  anos)
\end{itemize}

\section{Quais são as etapas para a realização desse serviço?}

Para retirada da CIP, são necessários dois procedimentos:

\begin{enumerate}
  \item  primeira participação no CRP Acolhe (encontro das/os  recém-inscritas/os com conselheiras/os para apresentação da  instituição e da legislação que norteia a atuação profissional da/o  psicóloga/o):

        \begin{enumerate}
          \item    \href{https://www.crpsp.org/inscrito/login}{agendamento, pelo    \emph{site}}, para participação no CRP Acolhe. A reunião pode    ocorrer de forma \emph{on-line} ou presencial;
          \item    participação na reunião agendada. A reunião tem duração de cerca de    2 horas.
        \end{enumerate}

  \item  Comparecimento na subsede de referência (confira os endereços  \href{https://www.crpsp.org/sede/index}{no \emph{site}}), mediante  \href{https://www.crpsp.org/agendamento}{agendamento prévio}, para  apresentar os documentos originais e coletar os dados  biométricos.
\end{enumerate}

Obs.: as etapas 1 e 2 poderão ser concluídas em qualquer ordem. Porém, aretirada da CIP somente será possível após a participação no CRP Acolhe.

\section{Quanto custa?}

\begin{itemize}
  \item  Inscrição: R\$~105,37 (valor para 2024);
  \item  emissão de 2ª~CIP: R\$~70,83 (valor para 2024)
\end{itemize}

\section{Quanto tempo leva?}

A emissão da CIP se dá no prazo de 30 dias úteis após a apresentação dosdocumentos e coleta dos dados biométricos. Para retirada, sãoobrigatórias a participação prévia no CRP Acolhe e a apresentação dedocumento de identidade com foto.

\section{Legislação relacionada}

\begin{itemize}
  \item  \href{https://atosoficiais.com.br/cfp/resolucao-administrativa-financeira-n-3-2007-institui-a-consolidacao-das-resolucoes-do-conselho-federal-de-psicologia?origin=instituicao&q=resolu\%C3\%A7\%C3\%A3o\%203/2007}{Resolução  CFP nº~03/2007} (Consolidação das Resoluções do Conselho Federal de  Psicologia);
  \item  \href{https://atosoficiais.com.br/cfp/resolucao-administrativa-financeira-n-20-2018-altera-o-manual-de-procedimentos-administrativo-e-financeiro-do-sistema-conselhos-de-psicologia-anexo-da-resolucao-cfp-n-202018-a-resolucao-cfp-n-03-2007-a-resolucao-cfp-n-16-2019-e-da-outras-providencias}{Resolução  CFP nº~20/2018}, que determina a revisão e ampliação do  \href{https://transparencia.cfp.org.br/wp-content/uploads/sites/7/2020/05/Manual-de-Procedimentos-Administrativos-e-Financeiro-CFP-RES-20.2018.pdf}{\emph{Manual  de procedimentos administrativos e financeiros}};
  \item  \href{https://atosoficiais.com.br/cfp/resolucao-administrativa-financeira-n-2-2021-dispoe-sobre-a-confeccao-a-expedicao-e-o-recolhimento-de-carteiras-de-identidade-profissional-das-os-psicologas-os-revoga-o-capitulo-vi-da-resolucao-cfp-no-03-2007-os-1o-2o-4o-5o-do-art-2o-da-resolucao-cfp-no-02-2007-1o-4o-5o-do-art-2o-e-1o-do-art-4o-da-resolucao-cfp-no-01-1984-e-da-outras-providencias}{Resolução  CFP nº~02/2021}, que dispõe sobre a CIP
\end{itemize}

\section{Serviços correlatos}

\begin{itemize}
  \item  \hyperref[]{Inscrição  profissional de pessoa física}
\end{itemize}

\chapter{\texorpdfstring{\emph{Alteração de inscrição provisória paradefinitiva}}{Alteração de inscrição provisória para definitiva}}


\section{O que é?}

Em até dois anos de inscrição provisória (realizada apenas com aapresentação do certificado de colação de grau), a/o profissional deveapresentar o diploma de formação, para que a inscrição sejadefinitiva.

A não apresentação do diploma no prazo acarretará o cancelamento dainscrição.

\section{Quem pode utilizar esse serviço?}

Profissionais graduadas/os em Psicologia, que tenham atuação no estadode São Paulo e que tenham realizado a inscrição no CRP~SP com a certidãode colação de grau.

\section{Como solicitar o serviço?}

A inscrição definitiva pode ser solicitada \emph{on-line}. Seránecessário realizar o\href{https://www.crpsp.org/pagina/view/311}{agendamento} no \emph{site}do CRP~SP, comparecendo à subsede de referência (confira os endereços\href{https://www.crpsp.org/sede/index}{no \emph{site}}) com o diplomaoriginal.

\section{Que informações ou documentos são necessários?}

\begin{itemize}
  \item  documento de identificação (RG, CNH etc.);
  \item  Cadastro de Pessoa Física (CPF);
  \item  diploma de formação em Psicologia (frente e verso);
  \item  CIP provisória
\end{itemize}

\section{Quais são as etapas para a realização desse serviço?}

\begin{enumerate}
  \item  Realizar a solicitação do serviço pelo \emph{site};
  \item  agendar dia e horário para apresentação do documento original.
\end{enumerate}

\section{Quanto custa?}

O processo de alteração de inscrição provisória para definitiva não temcusto. Porém, a CIP definitiva terá o mesmo custo de emissão da 2ª~viada CIP.

\section{Quanto tempo leva?}

Após a submissão do pedido no \emph{site}, o prazo é de 30 dias.

\section{Como ter acesso ao atual \emph{status} do serviço?}

No \emph{site} do CRP~SP, acessando a página da/o psicóloga/o com seu\emph{login} e senha, no campo\href{https://cfpsp.brctotal.com/crp06_servicosonline/pgsRequerimento/MeusRequerimentos.aspx}{Meusrequerimentos}.

\section{Legislação relacionada}

\begin{itemize}
  \item  Resolução CFP nº~03/2007 (Consolidação das Resoluções do Conselho  Federal de Psicologia);
  \item  \href{https://transparencia.cfp.org.br/wp-content/uploads/sites/7/2020/05/Manual-de-Procedimentos-Administrativos-e-Financeiro-CFP-RES-20.2018.pdf}{\href{https://atosoficiais.com.br/cfp/resolucao-administrativa-financeira-n-20-2018-altera-o-manual-de-procedimentos-administrativo-e-financeiro-do-sistema-conselhos-de-psicologia-anexo-da-resolucao-cfp-n-202018-a-resolucao-cfp-n-03-2007-a-resolucao-cfp-n-16-2019-e-da-outras-providencias}{Resolução  CFP nº~20/2018}}, que determina a revisão e ampliação do \emph{Manual  de procedimentos administrativos e financeiros}  (\href{https://transparencia.cfp.org.br/wp-content/uploads/sites/7/2020/05/Manual-de-Procedimentos-Administrativos-e-Financeiro-CFP-RES-20.2018.pdf}{versão  atualizada})
\end{itemize}

\section{Serviços correlatos}

\begin{itemize}
  \item  \hyperref[]{Inscrição  profissional de pessoa física};
  \item  \hyperref[]{entrega  da carteira de identidade profissional}
\end{itemize}

\chapter{\texorpdfstring{\emph{Prorrogação do prazo de entrega dodiploma}}{Prorrogação do prazo de entrega do diploma}}
\hyperref[]{Quais são as etapas para a realização desse serviço?}

\hyperref[]{Quanto custa?}

\hyperref[]{Quanto tempo leva?}

\hyperref[]{Como ter acesso ao atual \emph{status} do serviço?}

\hyperref[]{Legislação relacionada}

\hyperref[]{Serviços correlatos}

\section{O que é?}

É a extensão em seis meses do prazo de entrega do diploma, paraalteração da inscrição de provisória para definitiva.

\section{Quem pode utilizar esse serviço?}

Psicólogas/os que, tendo solicitado o diploma em tempo hábil, tenham odocumento retido ou não expedido pela entidade formadora devido adébitos em aberto.

\section{Como solicitar o serviço?}

Se não estiver em posse do diploma até o vencimento da inscriçãoprovisória, solicite a prorrogação pelo período de até seis meses,enviando o protocolo ou a declaração do pedido de emissão do diplomafornecido pela entidade formadora (IES) para a sua subsede.

\section{Que informações ou documentos são necessários?}

Protocolo ou declaração do pedido de emissão do diploma fornecido pelaentidade formadora (IES)

\section{Quais são as etapas para a realização desse serviço?}

\begin{enumerate}
  \item  A/o psicóloga/o solicita a prorrogação do prazo de entrega do diploma  antes do prazo de dois anos, contados da solicitação de inscrição  provisória;
  \item  o CRP~SP analisa a documentação enviada e concede a extensão do prazo.
\end{enumerate}

\subsection{Se a/o psicóloga/o apresenta o diploma no prazo de seis meses}

\begin{enumerate}
  \item  \hyperref[]{A  carteira de identidade profissional é entregue}.
\end{enumerate}

\subsection{Se a/o psicóloga/o não apresenta o diploma no prazo de seis meses}

\begin{enumerate}
  \item  O CRP~SP notifica a/o psicóloga/o do vencimento do prazo concedendo 30  dias para a entrega do documento ou nova e última prorrogação do prazo  de entrega do documento;
\end{enumerate}

\begin{enumerate}
  \item  a/o psicóloga/o encaminha nova solicitação de prorrogação, junto com o  protocolo de solicitação da entidade formadora;
  \item  se o diploma for entregue no prazo da segunda prorrogação, a  \hyperref[]{CIP  definitiva é expedida e entregue}. Caso contrário, o CRP~SP procede ao  cancelamento da inscrição provisória.
\end{enumerate}

\section{Quanto custa?}

Não tem custo.

\section{Quanto tempo leva?}

Após a submissão do pedido no \emph{site}, o prazo é de até 30 dias.

\section{Como ter acesso ao atual \emph{status} do serviço?}

No \emph{site} do CRP~SP, acessando a página da/o psicóloga/o com seu\emph{login} e senha, no campo\href{https://cfpsp.brctotal.com/crp06_servicosonline/pgsRequerimento/MeusRequerimentos.aspx}{Meusrequerimentos}.

\section{Legislação relacionada}

\begin{itemize}
  \item  \href{https://atosoficiais.com.br/cfp/resolucao-administrativa-financeira-n-3-2007-institui-a-consolidacao-das-resolucoes-do-conselho-federal-de-psicologia?origin=instituicao&q=resolu\%C3\%A7\%C3\%A3o\%203/2007}{Resolução  CFP nº~03/2007} (Consolidação das Resoluções do Conselho Federal de  Psicologia);
  \item  \href{https://transparencia.cfp.org.br/wp-content/uploads/sites/7/2020/05/Manual-de-Procedimentos-Administrativos-e-Financeiro-CFP-RES-20.2018.pdf}{\href{https://atosoficiais.com.br/cfp/resolucao-administrativa-financeira-n-20-2018-altera-o-manual-de-procedimentos-administrativo-e-financeiro-do-sistema-conselhos-de-psicologia-anexo-da-resolucao-cfp-n-202018-a-resolucao-cfp-n-03-2007-a-resolucao-cfp-n-16-2019-e-da-outras-providencias}{Resolução  CFP nº~20/2018}}, que determina a revisão e ampliação do  \href{https://transparencia.cfp.org.br/wp-content/uploads/sites/7/2020/05/Manual-de-Procedimentos-Administrativos-e-Financeiro-CFP-RES-20.2018.pdf}{\emph{Manual  de procedimentos administrativos e financeiros}}
\end{itemize}

\section{Serviços correlatos}

\begin{itemize}
  \item  \hyperref[]{Entrega  da carteira de identidade profissional};
  \item  \hyperref[]{alteração  da inscrição de provisória para definitiva}
\end{itemize}

\chapter{Emissão de 2ª~via de CIP}


\section{O que é?}

É a emissão de nova CIP devido a alteração de nome, inserção de registrode especialista, alteração de inscrição provisória para definitiva,roubo ou extravio do documento original.

\section{Quem pode utilizar esse serviço?}

Profissionais graduadas/os em Psicologia que atuem no estado de SãoPaulo e já tenham inscrição no CRP~SP ativa.

\section{Como solicitar o serviço?}

\href{https://www.crpsp.org/agendamento}{Agendamento prévio} paraapresentação dos documentos originais na subsede de referência (confiraos endereços \href{https://www.crpsp.org/sede/index}{no \emph{site}}).

\section{Que informações ou documentos são necessários?}

\begin{itemize}
  \item  Documento de identificação (RG, CNH etc.);
  \item  Cadastro de Pessoa Física (CPF);
  \item  diploma de formação em Psicologia;
  \item  deferimento de registro de especialista;
  \item  certidão de casamento ou de união estável;
  \item  boletim de ocorrência (se a CIP foi roubada ou extraviada)
\end{itemize}

\section{Quais são as etapas para a realização desse serviço?}

\begin{enumerate}
  \item  A/o psicóloga/o agenda o atendimento e paga a taxa de emissão da CIP;
\end{enumerate}

\begin{enumerate}
  \item  a/o psicóloga/o apresenta os documentos originais na subsede de  referência (confira os endereços  \href{https://www.crpsp.org/sede/index}{no \emph{site}}) e, se  necessário, faz a coleta de dados biométricos;
  \item  o CRP~SP emite a segunda via da CIP.
\end{enumerate}

\section{Quanto custa?}

\begin{itemize}
  \item  R\$~70,83 (valor de emissão da CIP em 2024)
\end{itemize}

\section{Quanto tempo leva?}

30 dias úteis a partir da data de solicitação presencial e coleta dosdados biométricos.

\section{Como ter acesso ao atual \emph{status} do serviço?}

{[}O CRP~SP comunica à/ao profissional que a CIP está disponível? Mandapelo correio?{]}

\section{Legislação relacionada}

\begin{itemize}
  \item  \href{https://atosoficiais.com.br/cfp/resolucao-administrativa-financeira-n-3-2007-institui-a-consolidacao-das-resolucoes-do-conselho-federal-de-psicologia?origin=instituicao&q=resolu\%C3\%A7\%C3\%A3o\%203/2007}{Resolução  CFP nº~03/2007} (Consolidação das Resoluções do Conselho Federal de  Psicologia);
  \item  \href{https://transparencia.cfp.org.br/wp-content/uploads/sites/7/2020/05/Manual-de-Procedimentos-Administrativos-e-Financeiro-CFP-RES-20.2018.pdf}{\href{https://atosoficiais.com.br/cfp/resolucao-administrativa-financeira-n-20-2018-altera-o-manual-de-procedimentos-administrativo-e-financeiro-do-sistema-conselhos-de-psicologia-anexo-da-resolucao-cfp-n-202018-a-resolucao-cfp-n-03-2007-a-resolucao-cfp-n-16-2019-e-da-outras-providencias}{Resolução  CFP nº~20/2018}}, que determina a revisão e ampliação do  \href{https://transparencia.cfp.org.br/wp-content/uploads/sites/7/2020/05/Manual-de-Procedimentos-Administrativos-e-Financeiro-CFP-RES-20.2018.pdf}{\emph{Manual  de procedimentos administrativos e financeiros}};
  \item  \href{https://atosoficiais.com.br/cfp/resolucao-administrativa-financeira-n-2-2021-dispoe-sobre-a-confeccao-a-expedicao-e-o-recolhimento-de-carteiras-de-identidade-profissional-das-os-psicologas-os-revoga-o-capitulo-vi-da-resolucao-cfp-no-03-2007-os-1o-2o-4o-5o-do-art-2o-da-resolucao-cfp-no-02-2007-1o-4o-5o-do-art-2o-e-1o-do-art-4o-da-resolucao-cfp-no-01-1984-e-da-outras-providencias}{Resolução  CFP nº~02/2021}, que trata da confecção, a expedição e o recolhimento  de carteiras de identidade profissional das/os psicólogas/os
\end{itemize}

\section{Serviços correlatos}

\begin{itemize}
  \item  \hyperref[]{Entrega  da carteira de identidade profissional (CIP)}
\end{itemize}

\chapter{\texorpdfstring{\emph{Solicitação de registro deespecialista}}{Solicitação de registro de especialista}}


\section{O que é?}

O registro de especialista atesta a experiência profissional em uma oumais\href{https://site.cfp.org.br/servicos/titulo-de-especialista/}{áreas deespecialidade} reconhecidas pelo Conselho Federal de Psicologia (CFP) enão constitui condição obrigatória para o exercício profissional.

O registro de especialista do CFP não é o mesmo que o título acadêmicoconferido a quem conclui uma pós-graduação \emph{lato sensu}:

\begin{itemize}
  \item  especialista é a/o profissional graduada/o que acumula notório saber  teórico-prático em uma área específica, reconhecida por seu conselho  de classe. Esse reconhecimento favorece a institucionalização da  profissão na sociedade brasileira e ajuda a demarcar suas áreas de  atuação;
  \item  especialização é o aprimoramento profissional obtido por meio de curso  de pós-graduação em uma área específica, que atenda às normativas e  aos critérios de funcionamento determinados pelo MEC ou outro órgão  com essa competência
\end{itemize}

Desse modo, há diferença entre ser psicóloga/o especialista epsicóloga/o com especialização: enquanto a/o primeira/o tem sua formaçãoe prática profissional reconhecidas/os por seu conselho regional, a/osegunda/o concluiu um curso de especialização.

\section{Quem pode utilizar esse serviço?}

Qualquer psicóloga/o que possua o CRP ativo, com no mínimo 2 (anos) deinscrição, que esteja adimplente com as anuidades dos exercíciosanteriores e que não esteja cumprindo penalidade resultante de processoético de multa, suspensão ou cassação.

\section{Como solicitar o serviço?}

Os pedidos são recebidos diretamente pelos\href{https://crpsp.org/sede/index}{\emph{e-mail}s institucionais dassubsedes}.

\section{Que informações ou documentos são necessários?}

\begin{itemize}
  \item  Certificado e histórico de curso de pós-graduação (reconhecido pelo  MEC) em alguma das  \href{https://site.cfp.org.br/servicos/titulo-de-especialista/}{áreas  definidas pelo CFP}; ou
  \item  documento que comprove a homologação do resultado final de concurso  para registro de especialista promovido pelo CFP;
  \item  comprovação de exercício profissional de no mínimo dois anos na  especialidade requerida (para cursos e provas concluídos a partir de 2  de fevereiro de 2023);
  \item  \hyperref[]{requerimento de  registro de especialista} preenchido e assinado
\end{itemize}

Mais informações estão disponíveis no\href{http://www.crpsp.org/pagina/view/51}{\emph{site} do CRP~SP}.

\section{Quais são as etapas para a realização desse serviço?}

\begin{enumerate}
  \item  A/o psicóloga/o solicita o registro de especialista via \emph{e-mail},  anexando o  \hyperref[]{requerimento de  registro de especialista} e cópias (escaneadas) dos documentos  necessários para cada situação;
\end{enumerate}

\begin{enumerate}
  \item  os documentos são analisados pela Comissão de Análise para Concessão  de Registro de Psicóloga/o Especialista (CARPE);
  \item  a Plenária do CRP~vota pelo deferimento ou indeferimento do registro.
\end{enumerate}

\section{Quanto custa?}

Gratuito. Será cobrada posteriormente taxa de emissão de nova CIP com oacréscimo de registro de especialista, caso o pedido seja deferido.

\section{Quanto tempo leva?}

60 dias a partir da data de protocolo do documento, conforme\href{https://atosoficiais.com.br/cfp/resolucao-do-exercicio-profissional-n-23-2022-institui-condicoes-para-concessao-e-registro-de-psicologa-e-psicologo-especialistas-reconhece-as-especialidades-da-psicologia-e-revoga-as-resolucoes-cfp-n-13-de-14-de-setembro-de-2007-n-3-de-5-de-fevereiro-de-2016-n-18-de-5-de-setembro-de-2019}{ResoluçãoCFP nº~23/2022}.

Caso a documentação esteja incompleta, a CARPE solicitará regularização,estendendo o período de análise final em 30 dias.

Em caso de indeferimento, é possível interpor recurso, devendo a/opsicóloga/o manifestar interesse em até 30 dias de ciência doindeferimento. Não há prazo para análise do pedido de recurso pelo CFP.

\section{Como ter acesso ao atual \emph{status} do serviço?}

O andamento do pedido pode ser consultado por \emph{e-mail}, telefone oupresencialmente. Para casos de recurso, as informações sãodisponibilizadas pelo CFP a partir da identificação do protocolo SEIvinculado ao pedido e comunicado à/ao psicóloga/o pelo CRP~SP.

\section{Legislação relacionada}

\begin{itemize}
  \item  \href{https://site.cfp.org.br/wp-content/uploads/2008/08/Resolucao_CFP_nx_013-2007.pdf}{Resolução  CFP nº~13/2007}, que institui a Consolidação das Resoluções relativas  ao Título Profissional de Especialista em Psicologia e dispõe sobre  normas e procedimentos para seu registro;
  \item  \href{https://atosoficiais.com.br/cfp/resolucao-do-exercicio-profissional-n-23-2022-institui-condicoes-para-concessao-e-registro-de-psicologa-e-psicologo-especialistas-reconhece-as-especialidades-da-psicologia-e-revoga-as-resolucoes-cfp-n-13-de-14-de-setembro-de-2007-n-3-de-5-de-fevereiro-de-2016-n-18-de-5-de-setembro-de-2019}{Resolução  CFP nº~23/2022}, que institui condições para concessão e registro de  psicóloga e psicólogo especialistas e reconhece as especialidades da  Psicologia
\end{itemize}

\section{Serviços correlatos}

\begin{itemize}
  \item  \hyperref[]{Entrega  da carteira de identidade profissional (CIP)}
\end{itemize}

\chapter{\texorpdfstring{\emph{Inscriçãosecundária}}{Inscrição secundária}}


\section{O que é?}

Inscrição realizada nos casos em que a/o psicóloga/o é inscrita/o em CRPde outra região, mas precisa desenvolver atividades profissionais em SãoPaulo por período superior a 90 (noventa) dias por ano. A/o profissionalpoderá solicitar inscrição secundária no CRP~SP, sem cobrança deanuidade.

\section{Quem pode utilizar esse serviço?}

Psicóloga/o inscrita/o em CRP de outra região que precise trabalhar, porperíodo superior a 90 (noventa) dias por ano, no estado de São Paulo.

\section{Como solicitar o serviço?}

Deve ser solicitado de forma \emph{on-line}, pelo\href{https://www.crpsp.org/pagina/view/338}{\emph{site} do CRP~SP}.

\section{Que informações ou documentos são necessários?}

\begin{itemize}
  \item  Documento de identificação (RG, CNH etc.);
  \item  Cadastro de Pessoa Física (CPF);
  \item  certidão de quitação eleitoral;
  \item  comprovante ou declaração de endereço;
  \item  diploma de formação em Psicologia (frente e verso) ou certidão de  colação de grau;
  \item  certidão de casamento ou de união estável;
  \item  certidão de quitação militar ou certificado de reservista (para homens  até 45 anos);
  \item  CIP do CRP de origem;
  \item  declaração com o motivo da solicitação
\end{itemize}

\section{Quais são as etapas para a realização desse serviço?}

\begin{enumerate}
  \item  Solicitação via \href{https://www.crpsp.org/agendamento}{\emph{site}  do CRP~SP};
  \item  comparecimento na subsede de referência (confira os endereços  \href{https://www.crpsp.org/sede/index}{no \emph{site}}), após  agendamento, para apresentação dos documentos originais e coleta dos  dados biométricos.
\end{enumerate}

\section{Quanto custa?}

\begin{itemize}
  \item  R\$~70,83 (valor de emissão da CIP em 2024; não há cobrança de  anuidade)
\end{itemize}

\section{Quanto tempo leva?}

30 dias após a solicitação no \emph{site}.

\section{Como ter acesso ao atual \emph{status} do serviço?}

No \emph{site} do CRP~SP, acessando a página do psicólogo com seu\emph{login} e senha, no campo\href{https://cfpsp.brctotal.com/crp06_servicosonline/pgsRequerimento/MeusRequerimentos.aspx}{Meusrequerimentos}.

\section{Legislação relacionada}

\begin{itemize}
  \item  \href{https://atosoficiais.com.br/cfp/resolucao-administrativa-financeira-n-3-2007-institui-a-consolidacao-das-resolucoes-do-conselho-federal-de-psicologia?origin=instituicao&q=resolu\%C3\%A7\%C3\%A3o\%203/2007}{Resolução  CFP nº~03/2007} (Consolidação das Resoluções do Conselho Federal de  Psicologia);
  \item  \href{https://transparencia.cfp.org.br/wp-content/uploads/sites/7/2020/05/Manual-de-Procedimentos-Administrativos-e-Financeiro-CFP-RES-20.2018.pdf}{\href{https://atosoficiais.com.br/cfp/resolucao-administrativa-financeira-n-20-2018-altera-o-manual-de-procedimentos-administrativo-e-financeiro-do-sistema-conselhos-de-psicologia-anexo-da-resolucao-cfp-n-202018-a-resolucao-cfp-n-03-2007-a-resolucao-cfp-n-16-2019-e-da-outras-providencias}{Resolução  CFP nº~20/2018}}, que determina a revisão e ampliação do  \href{https://transparencia.cfp.org.br/wp-content/uploads/sites/7/2020/05/Manual-de-Procedimentos-Administrativos-e-Financeiro-CFP-RES-20.2018.pdf}{\emph{Manual  de procedimentos administrativos e financeiros}}
\end{itemize}

\section{Serviços correlatos}

\begin{itemize}
  \item  \hyperref[]{Inscrição  profissional de pessoa física};
  \item  \hyperref[]{entrega  da carteira de identidade profissional (CIP)}
\end{itemize}

\chapter{\texorpdfstring{\emph{Transferência de registroprofissional}}{Transferência de registro profissional}}


\section{O que é?}

Transferência de registro profissional da/o psicóloga/o com inscrição emoutro CRP para o CRP~SP.

\section{Quem pode utilizar esse serviço?}

Profissionais inscritas/os em outra jurisdição que passaram a residir noestado de São Paulo.

\section{Como solicitar o serviço?}

A solicitação deve ser feita\href{https://www.crpsp.org/pagina/view/337}{no \emph{site} do CRP~SP}.

\section{Que informações ou documentos são necessários?}

\begin{itemize}
  \item  Documento de identificação (RG, CNH etc.);
  \item  Cadastro de Pessoa Física (CPF)
  \item  certidão de quitação eleitoral;
  \item  comprovante ou declaração de endereço;
  \item  diploma de formação em Psicologia (frente e verso) ou certidão de  colação de grau;
  \item  certidão de casamento ou de união estável;
  \item  certidão de quitação militar ou certificado de reservista (para homens  até 45 anos).
  \item  CIP do Conselho Regional de origem
\end{itemize}

\section{Quais são as etapas para a realização desse serviço?}

\begin{enumerate}
  \item  A/o psicóloga/o realiza a solicitação de transferência \emph{on-line},  pelo \href{https://www.crpsp.org/pagina/view/337}{\emph{site} do  CRP~SP};
  \item  o CRP~SP faz a análise da documentação;
  \item  são emitidos os boletos para pagamento da taxa de inscrição e da  anuidade;
  \item  após o pagamento das taxas, o registro é deferido e uma declaração  para atuação profissional é disponibilizada.
\end{enumerate}

\section{Quanto custa?}

\begin{itemize}
  \item  Taxa administrativa: R\$~70,83 (valor para 2024);
  \item  anuidade vigente proporcional (consulte a  \href{https://www.crpsp.org/pagina/view/206}{Tabela de Anuidades e  Taxas --- Pessoa Física})
\end{itemize}

\section{Quanto tempo leva?}

Até 30 dias.

Enquanto aguarda o processo de transferência, é permitido o exercício daprofissão fora da área de jurisdição de seu CRP de origem pelo prazo deaté 90 dias, conforme artigo 9, § 1º~da\href{https://atosoficiais.com.br/cfp/resolucao-administrativa-financeira-n-3-2007-institui-a-consolidacao-das-resolucoes-do-conselho-federal-de-psicologia?origin=instituicao&q=resolu\%C3\%A7\%C3\%A3o\%203/2007}{ResoluçãoCFP nº~03/2007}.

\section

No \emph{site} do CRP~SP, acessando a página do psicólogo com seu\emph{login} e senha, no campo\href{https://cfpsp.brctotal.com/crp06_servicosonline/pgsRequerimento/MeusRequerimentos.aspx}{Meusrequerimentos}.

\section{Legislação relacionada}

\begin{itemize}
  \item  \href{https://atosoficiais.com.br/cfp/resolucao-administrativa-financeira-n-3-2007-institui-a-consolidacao-das-resolucoes-do-conselho-federal-de-psicologia?origin=instituicao&q=resolu\%C3\%A7\%C3\%A3o\%203/2007}{Resolução  CFP nº~03/2007} (Consolidação das Resoluções do Conselho Federal de  Psicologia);
  \item  \href{https://transparencia.cfp.org.br/wp-content/uploads/sites/7/2020/05/Manual-de-Procedimentos-Administrativos-e-Financeiro-CFP-RES-20.2018.pdf}{\href{https://atosoficiais.com.br/cfp/resolucao-administrativa-financeira-n-20-2018-altera-o-manual-de-procedimentos-administrativo-e-financeiro-do-sistema-conselhos-de-psicologia-anexo-da-resolucao-cfp-n-202018-a-resolucao-cfp-n-03-2007-a-resolucao-cfp-n-16-2019-e-da-outras-providencias}{Resolução  CFP nº~20/2018}}, que determina a revisão e ampliação do  \href{https://transparencia.cfp.org.br/wp-content/uploads/sites/7/2020/05/Manual-de-Procedimentos-Administrativos-e-Financeiro-CFP-RES-20.2018.pdf}{\emph{Manual  de procedimentos administrativos e financeiros}}
\end{itemize}

\section{Serviços correlatos}

\begin{itemize}
  \item  \hyperref[]{Inscrição secundária}
\end{itemize}

\chapter{\texorpdfstring{\emph{Cancelamento de registroprofissional}}{Cancelamento de registro profissional}}

\section{O que é?}

Cancelamento de registro de psicóloga/o. O registro cancelado impede oexercício profissional.

\section{Quem pode utilizar esse serviço?}

Psicólogas/os inscritas/os e ativas/os, que não exercerão a profissãopelo período em que o registro estiver cancelado.

\section{Como solicitar o serviço?}

A solicitação de cancelamento de inscrição deverá ser enviada pelocorreio com aviso de recebimento (AR), acompanhada de carteira deidentidade profissional original e do formulário d\hl{e cancelamento,}disponível no \href{https://www.crpsp.org/pagina/view/312}{\emph{site}do CRP~SP}, \hl{devidamente assinado}.

\section{Que informações ou documentos são necessários?}

\begin{itemize}
  \item  \hyperref[]{Requerimento  de cancelamento de inscrição (pessoa física)} preenchido, datado e  assinado;
  \item  CIP original
\end{itemize}

\section{Quais são as etapas para a realização desse serviço?}

\begin{itemize}
  \item  Envio da documentação necessária via carta com aviso de recebimento. O  registro será considerado cancelado a partir da data de registro da  correspondência no correio; ou
  \item  entrega dos documentos necessários na subsede de referência (confira  os endereços \href{https://www.crpsp.org/sede/index}{no \emph{site}})
\end{itemize}

\section{Quanto custa?}

Sem custo.

Obs.: O cancelamento não desobriga a/o psicóloga/o do pagamento dedébitos anteriores à data de cancelamento.

\section{Quanto tempo leva?}

30 dias para os processos enviados via correio; imediato no caso desolicitação presencial.

\section{Como ter acesso ao atual \emph{status} do serviço?}

Pode-se consultar a inscrição no\href{https://cadastro.cfp.org.br/}{Cadastro Nacional de Profissionaisde Psicologia}.

\section{Legislação relacionada}

\begin{itemize}
  \item  \href{https://atosoficiais.com.br/cfp/resolucao-administrativa-financeira-n-3-2007-institui-a-consolidacao-das-resolucoes-do-conselho-federal-de-psicologia?origin=instituicao&q=resolu\%C3\%A7\%C3\%A3o\%203/2007}{Resolução  CFP nº~03/2007} (Consolidação das Resoluções do Conselho Federal de  Psicologia);
  \item  \href{https://transparencia.cfp.org.br/wp-content/uploads/sites/7/2020/05/Manual-de-Procedimentos-Administrativos-e-Financeiro-CFP-RES-20.2018.pdf}{\href{https://atosoficiais.com.br/cfp/resolucao-administrativa-financeira-n-20-2018-altera-o-manual-de-procedimentos-administrativo-e-financeiro-do-sistema-conselhos-de-psicologia-anexo-da-resolucao-cfp-n-202018-a-resolucao-cfp-n-03-2007-a-resolucao-cfp-n-16-2019-e-da-outras-providencias}{Resolução  CFP nº~20/2018}}, que determina a revisão e ampliação do  \href{https://transparencia.cfp.org.br/wp-content/uploads/sites/7/2020/05/Manual-de-Procedimentos-Administrativos-e-Financeiro-CFP-RES-20.2018.pdf}{\emph{Manual  de procedimentos administrativos e financeiros}}
\end{itemize}

\chapter{\texorpdfstring{\emph{Reativação de registroprofissional}}{Reativação de registro profissional}}


\section{O que é?}

É a reativação de registro profissional cancelado.

\section{Quem pode utilizar esse serviço?}

Profissionais que tenham cancelado seu registro profissional e quedesejem reativar.

\section{Como solicitar o serviço?}

Por meio do \emph{site} do CRP~SP, na guia\href{https://www.crpsp.org/pagina/view/314}{Atendimento/Pessoafísica}.

\section{Que informações ou documentos são necessários?}

\begin{itemize}
  \item  Documento de identificação (RG, CNH etc.);
  \item  Cadastro de Pessoa Física (CPF)
  \item  certidão de quitação eleitoral;
  \item  comprovante ou declaração de endereço;
  \item  diploma de formação em Psicologia (frente e verso);
  \item  certidão de casamento ou de união estável.
  \item  certidão de quitação militar ou certificado de reservista (para homens  até 45 anos)
\end{itemize}

\section{Quais são as etapas para a realização desse serviço?}

\begin{enumerate}
  \item  A/o psicóloga/o faz a  \href{https://www.crpsp.org/pagina/view/314}{solicitação de inscrição  prévia \emph{on-line}};
  \item  o CRP~SP faz a análise da documentação;
  \item  são emitidos os boletos para pagamento da taxa de inscrição e da  anuidade;
  \item  após o pagamento das taxas, o pedido é deferido e a declaração para  atuação profissional é enviada.
\end{enumerate}

\section{Quanto custa?}

\begin{itemize}
  \item  Taxa administrativa: R\$~105,37 (valor para 2024);
  \item  anuidade vigente proporcional (consulte a  \href{https://www.crpsp.org/pagina/view/206}{Tabela de Anuidades e  Taxas --- Pessoa Física})
\end{itemize}

\section{Quanto tempo leva?}

Até 30 dias.

\section{Como ter acesso ao atual \emph{status} do serviço?}

Ao se fazer a solicitação, é gerada uma senha de acesso. Com essa senhaé possível fazer o acompanhamento \emph{on-line} do \emph{status} dasolicitação, na página\href{https://cfpsp.brctotal.com/crp06_servicosonline/pgsRequerimento/Requerimentos.aspx}{Meusrequerimentos}.

\section{Legislação relacionada}

\begin{itemize}
  \item  \href{https://atosoficiais.com.br/cfp/resolucao-administrativa-financeira-n-3-2007-institui-a-consolidacao-das-resolucoes-do-conselho-federal-de-psicologia?origin=instituicao&q=resolu\%C3\%A7\%C3\%A3o\%203/2007}{Resolução  CFP nº~03/2007} (Consolidação das Resoluções do Conselho Federal de  Psicologia);
  \item  \href{https://transparencia.cfp.org.br/wp-content/uploads/sites/7/2020/05/Manual-de-Procedimentos-Administrativos-e-Financeiro-CFP-RES-20.2018.pdf}{\href{https://atosoficiais.com.br/cfp/resolucao-administrativa-financeira-n-20-2018-altera-o-manual-de-procedimentos-administrativo-e-financeiro-do-sistema-conselhos-de-psicologia-anexo-da-resolucao-cfp-n-202018-a-resolucao-cfp-n-03-2007-a-resolucao-cfp-n-16-2019-e-da-outras-providencias}{Resolução  CFP nº~20/2018}}, que determina a revisão e ampliação do  \href{https://transparencia.cfp.org.br/wp-content/uploads/sites/7/2020/05/Manual-de-Procedimentos-Administrativos-e-Financeiro-CFP-RES-20.2018.pdf}{\emph{Manual  de procedimentos administrativos e financeiros}}
\end{itemize}

\section{Serviços correlatos}

\begin{itemize}
  \item  \hyperref[]{Entrega  da carteira de identidade profissional (CIP)};
  \item  \hyperref[]{alteração  de inscrição provisória para definitiva};
  \item  \hyperref[]{inscrição secundária};
  \item  \hyperref[]{transferência de  registro profissional};
  \item  \hyperref[]{cancelamento de  registro profissional};
  \item  \hyperref[]{prorrogação  do prazo de entrega do diploma};
  \item  \hyperref[]{solicitação  de isenção e interrupção temporária de anuidades}
\end{itemize}

\chapter{\texorpdfstring{\emph{Ressarcimento depagamento}}{Ressarcimento de pagamento}}


\section{O que é?}

Devolução de valores pagos indevidamente ou em duplicidade, ou nos casosde isenção da anuidade ou cancelamento da inscrição (ressarcimentoproporcional).

\section{Quem pode utilizar esse serviço?}

Toda a categoria profissional de psicólogas/os.

\section{Como solicitar o serviço?}

Através do \emph{site}, na página\href{https://www.crpsp.org/pagina/view/27/anuidade-pessoa-fisica}{Anuidade--- pessoa física}.

\section{Que informações ou documentos são necessários?}

\begin{itemize}
  \item  Dados da/o psicóloga/o;
  \item  comprovantes de pagamento indevido e/ou em duplicidade;
  \item  dados bancários
\end{itemize}

\section{Quais são as etapas para a realização desse serviço?}

\begin{enumerate}
  \item  \href{https://cfpsp.brctotal.com/crp06_servicosonline/pgsRequerimento/SelecionaRequerimento.aspx}{Entrar  no \emph{site}} com \emph{login} e senha e solicitar o reembolso.
\end{enumerate}

\section{Quanto custa?}

Sem custo.

\section{Quanto tempo leva?}

30 dias úteis.

\section{Como ter acesso ao atual \emph{status} do serviço?}

No cadastro profissional no \emph{site}, acessados com \emph{login} esenha na página\href{https://cfpsp.brctotal.com/crp06_servicosonline/pgsRequerimento/Requerimentos.aspx}{Meusrequerimentos}.

\section{Legislação relacionada}

\begin{itemize}
  \item  \href{https://atosoficiais.com.br/cfp/resolucao-administrativa-financeira-n-3-2007-institui-a-consolidacao-das-resolucoes-do-conselho-federal-de-psicologia?origin=instituicao&q=resolu\%C3\%A7\%C3\%A3o\%203/2007}{Resolução  CFP nº~03/2007} (Consolidação das Resoluções do Conselho Federal de  Psicologia);
  \item  \href{https://transparencia.cfp.org.br/wp-content/uploads/sites/7/2020/05/Manual-de-Procedimentos-Administrativos-e-Financeiro-CFP-RES-20.2018.pdf}{\href{https://atosoficiais.com.br/cfp/resolucao-administrativa-financeira-n-20-2018-altera-o-manual-de-procedimentos-administrativo-e-financeiro-do-sistema-conselhos-de-psicologia-anexo-da-resolucao-cfp-n-202018-a-resolucao-cfp-n-03-2007-a-resolucao-cfp-n-16-2019-e-da-outras-providencias}{Resolução  CFP nº~20/2018}}, que determina a revisão e ampliação do  \href{https://transparencia.cfp.org.br/wp-content/uploads/sites/7/2020/05/Manual-de-Procedimentos-Administrativos-e-Financeiro-CFP-RES-20.2018.pdf}{\emph{Manual  de procedimentos administrativos e financeiros}}
\end{itemize}

\section{Serviços correlatos}

\begin{itemize}
  \item  \hyperref[]{Inscrição  profissional de pessoa física};
  \item  anuidade\end{itemize}

\chapter{\texorpdfstring{\emph{Solicitação de isenção e interrupçãotemporária deanuidades}}{Solicitação de isenção e interrupção temporária de anuidades}}

\section{O que é?}

Isenção da anuidade em razão de idade da/o profissional (a partir de 65anos) ou interrupção temporária das anuidades em razão de residência noexterior ou de afastamento por motivo de saúde.

\section{Quem pode utilizar esse serviço?}

\begin{itemize}
  \item  Psicólogas/os a partir de 65 anos, ainda que incompletos (conta-se o  ano de nascimento, somente);
  \item  psicólogas/os que tenham viagem ao exterior, com permanência superior  a seis meses de ausência do país, proporcional ao tempo em que a  psicóloga(o) estiver ausente do país, com requerimento limitado a 12  meses a contar da data de retorno ao país; e/ou
  \item  psicólogas/os em licença por motivo de saúde por período superior a 30  dias, proporcional ao tempo em que a psicóloga(o) estiver em  tratamento, com requerimento limitado a 12 meses a partir da alta  médica
\end{itemize}

\section{Como solicitar o serviço?}

\subsection{No caso de isenção poridade}

A isenção é concedida automaticamente.

\subsection{Nos demais casos}

Encaminhar a documentação comprobatória (cópias autenticadas) àpresidência do CRP~SP, pelo correio; ou entregar os documentos originaisna subsede de referência (confira os endereços\href{https://www.crpsp.org/sede/index}{no \emph{site}}).

\section{Que informações ou documentos são necessários?}

\subsection{No caso de isenção poridade}

Não é necessário enviar documentos.

\subsection{Nos demais casos}

Toda a documentação comprobatória da razão da isenção (laudo pericialemitido por profissional devidamente registrado no CRP ou no CRM,comprovantes de viagem ao exterior em que se ateste o período depermanência fora do país etc.).

\section{Quais são as etapas para a realização desse serviço?}

\begin{enumerate}
  \item  Encaminhar documentação comprobatória pelo correio ou entregar os  documentos originais na subsede de referência (confira os endereços  \href{https://www.crpsp.org/sede/index}{no \emph{site}}).
\end{enumerate}

\section{Quanto custa?}

Sem custo.

\section{Quanto tempo leva?}

A diretoria do Conselho Regional de Psicologia decidirá em dez dias,cabendo recurso ao Plenário, no prazo de 20 dias, em caso deindeferimento.

No caso de não deliberação no prazo de 30 dias da data do recebimento dopedido, a interrupção temporária será tida como aprovada.

\section{Como ter acesso ao atual \emph{status} do serviço?}

Uma vez recebidos e protocolados os documentos, é possível acompanhar oprocesso por meio do cadastro profissional no \emph{site}, acessado com\emph{login} e senha, na página\href{https://cfpsp.brctotal.com/crp06_servicosonline/pgsRequerimento/Requerimentos.aspx}{Meusrequerimentos}.

\section{Legislação relacionada}

\begin{itemize}
  \item  \href{https://atosoficiais.com.br/cfp/resolucao-administrativa-financeira-n-3-2007-institui-a-consolidacao-das-resolucoes-do-conselho-federal-de-psicologia?origin=instituicao&q=resolu\%C3\%A7\%C3\%A3o\%203/2007}{Resolução  CFP nº~03/2007} (Consolidação das Resoluções do Conselho Federal de  Psicologia);
  \item  \href{https://transparencia.cfp.org.br/wp-content/uploads/sites/7/2020/05/Manual-de-Procedimentos-Administrativos-e-Financeiro-CFP-RES-20.2018.pdf}{\href{https://atosoficiais.com.br/cfp/resolucao-administrativa-financeira-n-20-2018-altera-o-manual-de-procedimentos-administrativo-e-financeiro-do-sistema-conselhos-de-psicologia-anexo-da-resolucao-cfp-n-202018-a-resolucao-cfp-n-03-2007-a-resolucao-cfp-n-16-2019-e-da-outras-providencias}{Resolução  CFP nº~20/2018}}, que determina a revisão e ampliação do  \href{https://transparencia.cfp.org.br/wp-content/uploads/sites/7/2020/05/Manual-de-Procedimentos-Administrativos-e-Financeiro-CFP-RES-20.2018.pdf}{\emph{Manual  de procedimentos administrativos e financeiros}}
\end{itemize}

\section{Serviços correlatos}

\begin{itemize}
  \item  \hyperref[]{Inscrição  profissional de pessoa física}
\end{itemize}

\chapter{\texorpdfstring{\emph{Emissão dedeclarações}}{Emissão de declarações}}

\hyperref[]{O que são?}
\section{O que são?}

Documentos que apresentam a regularidade do profissional diante doConselho. Podem trazer informações referentes a:

\begin{itemize}
  \item  \hyperref[]{cancelamento de  registro profissional};
  \item  \hyperref[]{solicitação  de registro de especialista};
  \item  \hyperref[]{alteração,  substituição e/ou inclusão de Responsável Técnica/o};
  \item  atuação no exterior;
  \item  processos em andamento na Comissão de Ética
\end{itemize}

\section{Quem pode utilizar esse serviço?}

Toda a categoria de psicólogas/os.

\section{Como solicitar o serviço?}

Por meio dos serviços \emph{on-line} disponibilizados no \emph{site}, épossível \href{https://www.crpsp.org/pagina/view/326}{emitir declaraçãode regularidade simples} (situação da inscrição).

Caso o profissional necessite de uma declaração com dados adicionais,deverá solicitar em sua subsede de referência (confira os endereços\href{https://www.crpsp.org/sede/index}{no \emph{site}}).

\section{Que informações ou documentos são necessários?}

\subsection{Para emissão direto no\emph{site}}

\begin{itemize}
  \item  Acesso à página  \href{https://cfpsp.brctotal.com/crp06_servicosonline/pgsRequerimento/Requerimentos.aspx}{Meus  requerimentos} com \emph{login} e senha
\end{itemize}

\subsection{Demais casos}

\begin{itemize}
  \item  Carteira de identidade profissional (CIP) com dados da/o psicóloga/o
\end{itemize}

\section{Quais são as etapas para a realização desse serviço?}

\subsection{\texorpdfstring{No caso de emissão\emph{on-line}}{No caso de emissão on-line}}

\begin{enumerate}
  \item  Entrar no  \href{https://cfpsp.brctotal.com/crp06_servicosonline/pgsRequerimento/Requerimentos.aspx}{\emph{site}}  com \emph{login} e senha;
\end{enumerate}

\begin{enumerate}
  \item  solicitar a certidão.
\end{enumerate}

\subsection{Para emissão na subsede}

\begin{enumerate}
  \item  Acessar a opção ``Certidão com informações adicionais'' no  \href{https://www.crpsp.org/pagina/view/326}{\emph{site}};
  \item  escolher a subsede de referência (confira os endereços  \href{https://www.crpsp.org/sede/index}{no \emph{site}});
  \item  preencher o formulário de solicitação da certidão, justificando o  motivo do requerimento do documento.
\end{enumerate}

\section{Quanto custa?}

Sem custo.

\section{Quanto tempo leva?}

\subsection{\texorpdfstring{No caso de emissão\emph{on-line}}{No caso de emissão on-line}}

Imediata.

\subsection{No caso de solicitação com informações adicionais (viasubsede)}

Até cinco dias úteis.

\section{Como ter acesso ao atual \emph{status} do serviço?}

\subsection{\texorpdfstring{No caso de emissão\emph{on-line}}{No caso de emissão on-line}}

Imediata.

\subsection{No caso de solicitação viasubsede}

Acessar os canais de comunicação com o setor de atendimento da subsede(confira as informações \href{https://www.crpsp.org/sede/index}{no\emph{site}}).

\section{Legislação relacionada}

\begin{itemize}
  \item  Resolução CFP nº~03/2007, modificada pela  \href{https://atosoficiais.com.br/cfp/resolucao-administrativa-financeira-n-20-2018-altera-o-manual-de-procedimentos-administrativo-e-financeiro-do-sistema-conselhos-de-psicologia-anexo-da-resolucao-cfp-n-202018-a-resolucao-cfp-n-03-2007-a-resolucao-cfp-n-16-2019-e-da-outras-providencias}{Resolução  CFP nº~20/2018}
\end{itemize}

\section{Serviços correlatos}

\begin{itemize}
  \item  \hyperref[]{Inscrição  profissional de pessoa física}
\end{itemize}

\part{Serviços à pessoajurídica}

\chapter{\texorpdfstring{\emph{Inscrição profissional de pessoajurídica}}{Inscrição profissional de pessoa jurídica}}

\section{O que é?}

É o registro da pessoa jurídica (PJ) para prestação de serviços dePsicologia.

A PJ que presta serviços de Psicologia está obrigada a se registrar noConselho Regional de Psicologia em cuja jurisdição exerça suasatividades.

\section{Quem pode utilizar esse serviço?}

Pessoa jurídica que preste serviços de Psicologia.

\section{Como solicitar o serviço?}

De forma \emph{on-line}, por meio do \emph{link} na guia\href{https://www.crpsp.org/pagina/view/30}{Atendimento/Pessoa jurídica}(os procedimentos, formulários e modelos estão disponíveis na mesmapágina).

\section{Que informações ou documentos são necessários?}

\begin{itemize}
  \item  Contrato social ou estatuto e ata de constituição ou requerimento de  empresária/o ou documento similar;
  \item  \hyperref[]{requerimento  de inscrição (pessoa jurídica)} preenchido;
  \item  \hyperref[]{termo  de responsabilidade técnica (RT)} preenchido;
  \item  \hyperref[]{declaração  institucional para garantia do amplo e livre exercício profissional  por parte dos profissionais de Psicologia e Responsáveis Técnicos de  Psicologia de acordo com o Código de Ética Profissional da/o  Psicóloga/o e legislação cabível} preenchida
\end{itemize}

\section{Quais são as etapas para a realização desse serviço?}

\begin{enumerate}
  \item  Faz-se o requerimento de inscrição de PJ de forma \emph{on-line},  \href{https://www.crpsp.org/pagina/view/30}{pelo \emph{site} do CRP};
\end{enumerate}

\begin{enumerate}
  \item  os setores administrativo e técnico do CRP~SP analisam a documentação;
  \item  após concluída a análise, são emitidos os boletos para pagamento da  taxa de inscrição e da anuidade;
  \item  após o pagamento das taxas de inscrição e da anuidade, o requerimento  é deferido e o certificado de inscrição é enviado junto com ofício de  orientação técnica.
\end{enumerate}

\section{Quanto custa?}

Os valores de taxa de inscrição e anuidade são definidos de acordo com ocapital social da PJ. É possível consultar os valores\href{https://www.crpsp.org/pagina/view/207}{\emph{no site} do CRP~SP}.

\section{Quanto tempo leva?}

Até 20 dias a partir da solicitação para lançamento da inscrição nosistema + 20 dias para análise técnica.

\section{Como ter acesso ao atual \emph{status} do serviço?}

Ao se fazer a solicitação, é gerada uma senha de acesso. Com essa senhaé possível fazer o acompanhamento \emph{on-line} do \emph{status} dasolicitação.

\section{Legislação relacionada}

\begin{itemize}
  \item  \href{https://transparencia.cfp.org.br/wp-content/uploads/sites/7/2020/05/Manual-de-Procedimentos-Administrativos-e-Financeiro-CFP-RES-20.2018.pdf}{\href{https://atosoficiais.com.br/cfp/resolucao-administrativa-financeira-n-20-2018-altera-o-manual-de-procedimentos-administrativo-e-financeiro-do-sistema-conselhos-de-psicologia-anexo-da-resolucao-cfp-n-202018-a-resolucao-cfp-n-03-2007-a-resolucao-cfp-n-16-2019-e-da-outras-providencias}{Resolução  CFP nº~20/2018}}, que determina a revisão e ampliação do  \href{https://transparencia.cfp.org.br/wp-content/uploads/sites/7/2020/05/Manual-de-Procedimentos-Administrativos-e-Financeiro-CFP-RES-20.2018.pdf}{\emph{Manual  de procedimentos administrativos e financeiros}};
  \item  \href{https://atosoficiais.com.br/cfp/resolucao-administrativa-financeira-n-16-2019-altera-o-manual-de-procedimentos-administrativo-e-financeiro-do-sistema-conselhos-de-psicologia-anexo-da-resolucao-cfp-n-202018-a-resolucao-cfp-n-03-2007-a-resolucao-cfp-n-16-2019-e-da-outras-providencias}{Resolução  CFP nº~16/2019}, que normatiza o registro e o cadastro de pessoas  jurídicas;
  \item  \href{https://atosoficiais.com.br/cfp/resolucao-administrativa-financeira-n-8-2023-altera-o-manual-de-procedimentos-administrativo-e-financeiro-do-sistema-conselhos-de-psicologia-anexo-da-resolucao-cfp-n-202018-a-resolucao-cfp-n-03-2007-a-resolucao-cfp-n-16-2019-e-da-outras-providencias?origin=instituicao}{Resolução  CFP nº~08/2023}, que altera o \emph{Manual de procedimentos  administrativos e financeiros}
\end{itemize}

\section{Serviços correlatos}

\begin{itemize}
  \item  \hyperref[]{Alteração contratual};
  \item  \hyperref[]{renovação do  certificado de licença};
  \item  \hyperref[]{baixa de  responsabilidade técnica};
  \item  \hyperref[]{cancelamento de inscrição};
  \item  \hyperref[]{alteração,  substituição e/ou inclusão de Responsável Técnica/o};
  \item  \hyperref[]{solicitação  de isenção de anuidade}
\end{itemize}

\chapter{\texorpdfstring{\emph{Renovação do certificado delicença}}{Renovação do certificado de licença}}


\section{O que é?}

Renovação do certificado emitido pelo CRP~SP que autoriza a pessoajurídica (PJ) a prestar serviços em Psicologia, válido por 3 anos.

\section{Quem pode utilizar esse serviço?}

Pessoas jurídicas registradas ou cadastradas no CRP.

\section{Como solicitar o serviço?}

Por \emph{e-mail}. Os procedimentos e os formulários estão disponíveisno \emph{site} do CRP~SP, na guia \href{https://www.crpsp.org/pagina/view/30}{Atendimento/Pessoajurídica}.

\section{Que informações ou documentos são necessários?}

\begin{itemize}
  \item  \hyperref[]{Termo  de responsabilidade técnica e autorização para envio de comunicações  eletrônicas} preenchidos;
  \item  \hyperref[]{solicitação  de renovação de certificado} preenchida e assinada pelas/os sócias/os  administradoras/es;
  \item  \hyperref[]{declaração  institucional para garantia do amplo e livre exercício profissional}  preenchida em papel timbrado da instituição pública ou privada;
  \item  última alteração contratual, se houver
\end{itemize}

\section{Quais são as etapas para a realização desse serviço?}

\begin{itemize}
  \item  Envio da documentação necessária por \emph{e-mail};
  \item  após verificação da documentação encaminhada, o CRP~SP  lançará/atualizará a informação no sistema e emitirá o boleto para  pagamento da taxa de renovação;
  \item  Após confirmação do pagamento da taxa, é enviado novo certificado,  junto com ofício de orientação técnica, por \emph{e-mail}
\end{itemize}

\section{Quanto custa?}

A taxa de renovação é calculada de acordo com o capital social da PJ. Atabela com os valores está disponível no\href{https://www.crpsp.org/pagina/view/207}{\emph{site} do CRP~SP}.

\section{Quanto tempo leva?}

15 dias úteis.

\section{Como ter acesso ao atual \emph{status} do serviço?}

Aguardar o período previsto para o serviço. Consultas podem ser feitaspor \emph{e-mail.}

\section{Legislação relacionada}

\begin{itemize}
  \item  \href{https://atosoficiais.com.br/cfp/resolucao-administrativa-financeira-n-16-2019-altera-o-manual-de-procedimentos-administrativo-e-financeiro-do-sistema-conselhos-de-psicologia-anexo-da-resolucao-cfp-n-202018-a-resolucao-cfp-n-03-2007-a-resolucao-cfp-n-16-2019-e-da-outras-providencias}{Resolução  CFP nº~16/2019}, que normatiza o registro e o cadastro de pessoas  jurídicas
\end{itemize}

\section{Serviços correlatos}

\begin{itemize}
  \item  \hyperref[]{Inscrição  profissional de pessoa jurídica}
\end{itemize}

\chapter{\texorpdfstring{\emph{Alteraçãocontratual}}{Alteração contratual}}


\section{O que é?}

Qualquer mudança (nome, endereço, capital social, objeto socialetc.) em cláusula de contrato social, estatuto, requerimento deempresário ou documento similar.

\section{Quem pode utilizar esse serviço?}

Pessoas jurídicas (PJ) registradas ou cadastradas no CRP~SP.

\section{Como solicitar o serviço?}

Por \emph{e-mail}.

Os procedimentos e os formulários estão disponíveis no\href{https://www.crpsp.org/pagina/view/320}{\emph{site} do CRP~SP}.

\section{Que informações ou documentos são necessários?}

\begin{itemize}
  \item  \hyperref[]{Solicitação  de cancelamento/alteração de registro de pessoa jurídica} preenchida;
  \item  \hyperref[]{termo  de responsabilidade técnica e autorização para envio de comunicações  eletrônicas} preenchidos;
  \item  \hyperref[]{declaração  institucional para garantia do amplo e livre exercício profissional  por parte das/os profissionais de Psicologia e Responsáveis  Técnicas/os de Psicologia de acordo com o Código de Ética Profissional  da/o Psicóloga/o e legislação cabível} preenchida em papel timbrado da  instituição pública ou privada;
  \item  cópia do documento de alteração contratual ou de documento que  comprove a alteração
\end{itemize}

\section{Quais são as etapas para a realização desse serviço?}

\begin{enumerate}
  \item  A pessoa responsável pela PJ deve encaminhar os  \href{https://www.crpsp.org/pagina/view/320}{documentos necessários}  por \emph{e-mail} ao CRP~SP, para análise;
\end{enumerate}

\begin{enumerate}
  \item  após verificação da documentação enviada, é feita a atualização  cadastral no sistema;
  \item  a certidão de alteração contratual é deferida e enviada por  \emph{e-mail}.
\end{enumerate}

\section{Quanto custa?}

Sem custo.

\section{Quanto tempo leva?}

15 dias úteis.

\section{Como ter acesso ao atual \emph{status} do serviço?}

Aguardar o período previsto para o serviço. Consultas podem ser feitaspor \emph{e-mail.}

\section{Legislação relacionada}

\begin{itemize}
  \item  \href{https://atosoficiais.com.br/cfp/resolucao-administrativa-financeira-n-16-2019-altera-o-manual-de-procedimentos-administrativo-e-financeiro-do-sistema-conselhos-de-psicologia-anexo-da-resolucao-cfp-n-202018-a-resolucao-cfp-n-03-2007-a-resolucao-cfp-n-16-2019-e-da-outras-providencias}{Resolução  CFP nº~16/2019}, que normatiza o registro e o cadastro de pessoas  jurídicas
\end{itemize}

\section{Serviços correlatos}

\begin{itemize}
  \item  \hyperref[]{Inscrição  profissional de pessoa jurídica}
\end{itemize}

\chapter{\texorpdfstring{\emph{Cancelamento deinscrição}}{Cancelamento de inscrição}}


\section{O que é?}

Cancelamento da inscrição de pessoa jurídica (PJ) junto ao CRP~SP, emvirtude de encerramento das atividades ou por exclusão dos serviços dePsicologia.

\section{Quem pode utilizar esse serviço?}

Pessoas jurídicas registradas ou cadastradas no CRP que encerrem suasatividades ou deixem de prestar serviços de Psicologia.

\section{Como solicitar o serviço?}

Por \emph{e-mail}. Os procedimentos e os formulários estão disponíveisno \emph{site} do CRP~SP, na guia\href{https://www.crpsp.org/pagina/view/30}{Atendimento/Pessoajurídica}.

\section{Que informações ou documentos são necessários?}

\subsection{Se cancelamento se der por encerramento das atividades daPJ}

\begin{itemize}
  \item  Distrato social ou documento equivalente;
  \item  \hyperref[]{solicitação  de cancelamento/alteração de registro de pessoa jurídica} preenchida;
  \item  \hyperref[]{formulário  de declaração de desligamento de Responsável Técnica/o} preenchido;
  \item  cópia da baixa do CNPJ ou da baixa na prefeitura
\end{itemize}

\subsection{No caso de cancelamento por exclusão dos serviços dePsicologia pelaPJ}

\begin{itemize}
  \item  Alteração contratual;
  \item  \hyperref[]{solicitação  de cancelamento/alteração de registro de pessoa jurídica} preenchida;
  \item  \hyperref[]{formulário  de declaração de desligamento de Responsável Técnica/o} preenchido;
  \item  \hyperref[]{formulário  de informações sobre as atividades da pessoa jurídica para análise do  pedido de cancelamento} preenchido
\end{itemize}

\section{Quais são as etapas para a realização desse serviço?}

\begin{enumerate}
  \item  Envio da documentação necessária por \emph{e-mail};
  \item  após verificação da documentação encaminhada, o CRP~SP lança a  informação no sistema;
  \item  o pedido é deferido e a certidão de cancelamento é enviada por  \emph{e-mail}.
\end{enumerate}

\section{Quanto custa?}

Sem custo.

\section{Quanto tempo leva?}

15 dias úteis.

\section{Como ter acesso ao atual \emph{status} do serviço?}

Aguardar o período previsto para o serviço. Consultas podem ser feitaspor \emph{e-mail.}

\section{Legislação relacionada}

\begin{itemize}
  \item  \href{https://transparencia.cfp.org.br/wp-content/uploads/sites/7/2020/05/Manual-de-Procedimentos-Administrativos-e-Financeiro-CFP-RES-20.2018.pdf}{\href{https://atosoficiais.com.br/cfp/resolucao-administrativa-financeira-n-20-2018-altera-o-manual-de-procedimentos-administrativo-e-financeiro-do-sistema-conselhos-de-psicologia-anexo-da-resolucao-cfp-n-202018-a-resolucao-cfp-n-03-2007-a-resolucao-cfp-n-16-2019-e-da-outras-providencias}{Resolução  CFP nº~20/2018}}, que determina a revisão e ampliação do  \href{https://transparencia.cfp.org.br/wp-content/uploads/sites/7/2020/05/Manual-de-Procedimentos-Administrativos-e-Financeiro-CFP-RES-20.2018.pdf}{\emph{Manual  de procedimentos administrativos e financeiros}};
  \item  \href{https://atosoficiais.com.br/cfp/resolucao-administrativa-financeira-n-16-2019-altera-o-manual-de-procedimentos-administrativo-e-financeiro-do-sistema-conselhos-de-psicologia-anexo-da-resolucao-cfp-n-202018-a-resolucao-cfp-n-03-2007-a-resolucao-cfp-n-16-2019-e-da-outras-providencias}{Resolução  CFP nº~16/2019}, que normatiza o registro e o cadastro de pessoas  jurídicas
\end{itemize}

\section{Serviços correlatos}

\begin{itemize}
  \item  \hyperref[]{Inscrição  profissional de pessoa jurídica}
\end{itemize}

\chapter{\texorpdfstring{\emph{Solicitação de isenção deanuidade}}{Solicitação de isenção de anuidade}}


\section{O que é?}

Solicitação de isenção da anuidade para pessoa jurídica (PJ), em razãode inatividade ou de mudança de estado.

\section{Quem pode utilizar esse serviço?}

Pessoas jurídicas registradas que estejam inativas ou que sejamtransferidas para outro estado.

\section{Como solicitar o serviço?}

Por correio, com aviso de recebimento (AR). Os procedimentos e osformulários estão disponíveis no \emph{site} do CRP~SP, na guia\href{https://www.crpsp.org/pagina/view/30}{Atendimento/Pessoajurídica}.

\section{Que informações ou documentos são necessários?}

\subsection{No caso de pedido porinatividade:}

\begin{itemize}
  \item  \hyperref[]{Solicitação  de isenção dos débitos de pessoa jurídica por alteração contratual}  preenchida e assinada;
  \item  declaração de inatividade emitida pela Receita Federal do Brasil (para  os casos de inatividade até 2015); ou
  \item  declaração de débitos e créditos tributários federais (DCTF) (para os  casos de inatividade a partir de 2016)
\end{itemize}

\subsection{No caso de pedido por mudança deestado:}

Envio por correio, com AR, dos documentos abaixo listados:

\begin{itemize}
  \item  \hyperref[]{solicitação  de cancelamento/alteração de registro de pessoa jurídica} preenchida;
  \item  \hyperref[]{solicitação  de isenção dos débitos de pessoa jurídica por alteração contratual}  preenchida;
  \item  alteração contratual (no caso de mudança para outro estado)
\end{itemize}

\section{Quais são as etapas para a realização desse serviço?}

\begin{itemize}
  \item  Enviar, por correio, a documentação necessária;
  \item  o CRP~SP analisa a documentação encaminhada;
  \item  é enviado, por \emph{e-mail}, ofício informando do deferimento da  solicitação
\end{itemize}

\section{Quanto custa?}

Sem custo.

\section{Quanto tempo leva?}

De 15 a 30 dias.

\section{Como ter acesso ao atual \emph{status} do serviço?}

Aguardar o período previsto para o serviço. Consultas podem ser feitaspor \emph{e-mail.}

\section{Legislação relacionada}

\begin{itemize}
  \item  \href{https://transparencia.cfp.org.br/wp-content/uploads/sites/7/2020/05/Manual-de-Procedimentos-Administrativos-e-Financeiro-CFP-RES-20.2018.pdf}{\href{https://atosoficiais.com.br/cfp/resolucao-administrativa-financeira-n-20-2018-altera-o-manual-de-procedimentos-administrativo-e-financeiro-do-sistema-conselhos-de-psicologia-anexo-da-resolucao-cfp-n-202018-a-resolucao-cfp-n-03-2007-a-resolucao-cfp-n-16-2019-e-da-outras-providencias}{Resolução  CFP nº~20/2018}}, que determina a revisão e ampliação do  \href{https://transparencia.cfp.org.br/wp-content/uploads/sites/7/2020/05/Manual-de-Procedimentos-Administrativos-e-Financeiro-CFP-RES-20.2018.pdf}{\emph{Manual  de procedimentos administrativos e financeiros}}
\end{itemize}

\section{Serviços correlatos}

\begin{itemize}
  \item  \hyperref[]{Inscrição  profissional de pessoa jurídica}
\end{itemize}

\chapter{\texorpdfstring{\emph{Alteração, substituição e/ou inclusãode ResponsávelTécnica/o}}{Alteração, substituição e/ou inclusão de Responsável Técnica/o}}


\section{O que é?}

Atualização do cadastro da pessoa jurídica (PJ), nos casos em que houveralteração, substituição e/ou inclusão de Responsável Técnica/o pelosserviços de Psicologia prestados pela empresa.

\section{Quem pode utilizar esse serviço?}

Pessoas jurídicas registradas ou cadastradas no CRP~SP.

\section{Como solicitar o serviço?}

Por \emph{e-mail}. Os procedimentos e formulários estão disponíveis no\href{https://www.crpsp.org/pagina/view/322}{\emph{site} do CRP~SP}.

\section{Que informações ou documentos são necessários?}

\begin{itemize}
  \item  \hyperref[]{Termo  de responsabilidade técnica (RT) e autorização para envio de  comunicações eletrônicas} preenchidos;
  \item  \hyperref[]{declaração  institucional para garantia do amplo e livre exercício profissional  por parte dos profissionais de Psicologia e Responsáveis Técnicos de  Psicologia de acordo com o Código de Ética Profissional da/o  Psicóloga/o e legislação cabível} preenchida em papel timbrado da  instituição pública ou privada;
  \item  carta da empresa informando a alteração, inclusão ou substituição da/o  Responsável Técnica/o;
  \item  última alteração contratual, se houver;
  \item  certificado de outro conselho de classe, se possuir
\end{itemize}

\section{Quais são as etapas para a realização desse serviço?}

\begin{enumerate}
  \item  Envio da documentação necessária por \emph{e-mail};
  \item  após verificação da documentação encaminhada, o CRP~SP  lançará/atualizará a informação no sistema;
  \item  novo certificado de inscrição é enviado, junto com ofício de  orientação técnica.
\end{enumerate}

\section{Quanto custa?}

Sem custo.

\section{Quanto tempo leva?}

De 15 a 30 dias.

\section{Como ter acesso ao atual \emph{status} do serviço?}

Aguardar o período previsto para o serviço. Consultas podem ser feitaspor \emph{e-mail}.

\section{Legislação relacionada}

\begin{itemize}
  \item  \href{https://atosoficiais.com.br/cfp/resolucao-administrativa-financeira-n-16-2019-altera-o-manual-de-procedimentos-administrativo-e-financeiro-do-sistema-conselhos-de-psicologia-anexo-da-resolucao-cfp-n-202018-a-resolucao-cfp-n-03-2007-a-resolucao-cfp-n-16-2019-e-da-outras-providencias}{Resolução  CFP nº~16/2019}, que normatiza o registro e o cadastro de pessoas  jurídicas
\end{itemize}

\section{Serviços correlatos}

\begin{itemize}
  \item  \hyperref[]{Inscrição  profissional de pessoa jurídica}
\end{itemize}

\chapter{\texorpdfstring{\emph{Baixa de responsabilidadetécnica}}{Baixa de responsabilidade técnica}}


\section{O que é?}

Desligamento ou afastamento da/o psicóloga/o que exerce a função deResponsável Técnica/o da pessoa jurídica (PJ).

\section{Quem pode utilizar esse serviço?}

Pessoas jurídicas ou a/o psicóloga/o RT da empresa.

\section{Como solicitar o serviço?}

Por \emph{e-mail}. Os procedimentos e os formulários estão disponíveisno \emph{site} do CRP~SP, na guia\href{https://www.crpsp.org/pagina/view/30}{Atendimento/Pessoajurídica}.

\section{Que informações ou documentos são necessários?}

\subsection{Psicóloga/o RT:}

\begin{itemize}
  \item  \hyperref[]{Formulário  de declaração de desligamento de Responsável Técnica/o} preenchido
\end{itemize}

\subsection{Pessoa jurídica:}

\begin{itemize}
  \item  Comunicar, dentro de 30 dias, a/o nova/o Responsável Técnica/o
\end{itemize}

É vedada a execução de serviços de Psicologia enquanto não houver adevida substituição.

\section{Quais são as etapas para a realização desse serviço?}

\begin{enumerate}
  \item  Envio da documentação necessária por \emph{e-mail};
  \item  após verificação da documentação encaminhada, o CRP~SP  lançará/atualizará a informação no sistema;
  \item  a confirmação de retirada de RT é enviada para a/o  profissional.
\end{enumerate}

\section{Quanto custa?}

Sem custo.

\section{Quanto tempo leva?}

15 dias úteis.

\section{Como ter acesso ao atual \emph{status} do serviço?}

Aguardar o período previsto para o serviço. Consultas podem ser feitaspor \emph{e-mail.}

\section{Legislação relacionada}

\begin{itemize}
  \item  \href{https://atosoficiais.com.br/cfp/resolucao-administrativa-financeira-n-16-2019-altera-o-manual-de-procedimentos-administrativo-e-financeiro-do-sistema-conselhos-de-psicologia-anexo-da-resolucao-cfp-n-202018-a-resolucao-cfp-n-03-2007-a-resolucao-cfp-n-16-2019-e-da-outras-providencias}{Resolução  CFP nº~16/2019}, que normatiza o registro e o cadastro de pessoas  jurídicas
\end{itemize}

\section{Serviços correlatos}

\begin{itemize}
  \item  \hyperref[]{Inscrição  profissional de pessoa jurídica};
  \item  \hyperref[]{alteração,  substituição e/ou inclusão de Responsável Técnica/o}
\end{itemize}

\section{Serviços de divulgação, orientação efiscalização}

{[}considerar a elaboração de página de abertura{]}

\chapter{\texorpdfstring{\emph{Acesso às publicaçõesorientativas}}{Acesso às publicações orientativas}}
\hyperref[]{Quais são as etapas para realização desse serviço?}

\hyperref[]{Quanto custa?}

\hyperref[]{Quanto tempo leva?}

\hyperref[]{Como ter acesso ao atual \emph{status} do serviço?}

\hyperref[]{Legislação relacionada}

\section{O que é?}

Coleção de documentos contendo orientações sobre legislaçãoprofissional, direcionados principalmente à categoria de psicólogas/os.Um exemplo de documento é a série\href{https://www.crpsp.org/impresso/index?categoria=9}{CRP~SP Orienta}.

\section{Quem pode utilizar esse serviço?}

Os materiais produzidos falam sobre a atuação ética e técnica deprofissionais de Psicologia. Podem ser destinados à categoria, ainstituições específicas ou às/aos usuárias/os do serviço.

\section{Como solicitar o serviço?}

Para acesso a orientações \emph{on-line}, basta acessar o\href{http://www.crpsp.org}{\emph{site} do CRP~SP} e clicar em``Orientação'' ou em ``Publicações''.

Para publicações impressas, é preciso verificar disponibilidade nas\href{https://www.crpsp.org/sede/index}{subsedes} do CRP~SP.

Para sugerir a confecção de novos materiais orientativos, envie um\emph{e-mail} para\href{mailto:orientacao@crpsp.org.br}{\nolinkurl{orientacao@crpsp.org.br}}indicando o conteúdo e sua importância. A Comissão de Orientação eFiscalização (COF) analisará o pedido e discutirá internamente suaelaboração.

\section{Que informações ou documentos são necessários?}

Para acessar o material já produzido não é necessária nenhumainformação.

\section{Quais são as etapas para a realização desse serviço?}

\subsection{Documentos já publicados}

Podem ser acessados no \href{http://www.crpsp.org}{\emph{site} doCRP~SP}, com destaque para as abas ``Orientação'' e ``Publicações'' do\emph{menu}.

\subsection{Novas publicações}

É preciso encaminhar por \emph{e-mail} o pedido e justificar suarelevância. O pedido será analisado pela COF, que fará consultas aoutros órgãos do CRP~SP e redigirá uma minuta de publicação paraaprovação pela Plenária. Uma vez aprovado, o material é encaminhado àComissão de Comunicação para criação de \emph{layout} e arte epublicação do documento.

\section{Quanto custa?}

Gratuito.

\section{Quanto tempo leva?}

Acesso imediato ao material já produzido. Não há prazo definido para aelaboração de novos conteúdos.

\section{Como ter acesso ao atual \emph{status} do serviço?}

Não se aplica.

\section{Legislação relacionada}

\begin{itemize}
  \item  \href{http://www.planalto.gov.br/ccivil_03/leis/l5766.htm}{Lei Federal  nº~5.766/1971}, que cria o Conselho Federal e os Conselhos Regionais  de Psicologia;
  \item  \href{https://atosoficiais.com.br/cfp/resolucao-do-exercicio-profissional-n-10-2005-aprova-o-codigo-de-etica-profissional-do-psicologo}{Resolução  CFP nº~10/2005}, que institui o  \href{https://www.crpsp.org/uploads/pagina/289379/2j9LIMPLJ9jFr5YrK57HmAiBWjVMxdbe.pdf}{Código  de Ética Profissional da/o Psicóloga/o}
\end{itemize}

\chapter{\texorpdfstring{\emph{Orientações sobre o Código de Ética elegislaçãoprofissional}}{Orientações sobre o Código de Ética e legislação profissional}}
\hyperref[]{Quais são as etapas para realização desse serviço?}

\hyperref[]{Quanto custa?}

\hyperref[]{Quanto tempo leva?}

\hyperref[]{Legislação relacionada}

\hyperref[]{Serviçoscorrelatos}

\section{O que é?}

A Comissão de Orientação e Fiscalização (COF) do CRP~SP, descentralizadanas 11 subsedes, oferece à categoria e à sociedade orientação sobreética e atuação profissional e temas relacionados à legislação doSistema Conselhos de Psicologia.

\section{Quem pode utilizar esse serviço?}

O acesso é disponibilizado à categoria e à sociedade, não serestringindo apenas às/aos cidadãs/ãos brasileiras/os.

\section{Como solicitar o serviço?}

O serviço é disponibilizado por telefone ou \emph{e-mail} e,presencialmente, pelos\href{http://www.crpsp.org/noticia/view/2460/contatos-das-subsedes-do-crp-sp}{contatosdiretos das subsedes}, e direcionado, de acordo com a demanda e omunicípio, para o território de referência.

O horário de atendimento (presencial e telefônico) da COF é das 10 às 16horas, de segunda a sexta-feira. O atendimento é feito em todas as 11subsedes do CRP~SP. Para atendimento presencial, recomendamosentrar em contato com a subsede antecipadamente.

\section{Que informações ou documentos são necessários?}

\subsection{Se o pedido é feito em umasubsede}

\begin{itemize}
  \item  A pessoa que solicita a informação deve se identificar, e os dados são  enviados ao Tribunal de Contas da União para controle (não é gerado  processo administrativo ou legal);
  \item  caso a/o solicitante seja psicóloga/o, também deverá informar o número  de inscrição no CRP
\end{itemize}

\chapter{\texorpdfstring{Se o pedido é feito por\emph{e-mail}}{Se o pedido é feito por e-mail}}

\begin{itemize}
  \item  Deve-se informar um \emph{e-mail} válido para envio da resposta
\end{itemize}

\section{Quais são as etapas para a realização desse serviço?}

\chapter{\texorpdfstring{Para orientações solicitadas via\emph{e-mail} dasubsede}{Para orientações solicitadas via e-mail da subsede}}

\begin{enumerate}
  \item  O pedido de orientação é feito por \emph{e-mail};
  \item  A resposta é enviada (também por \emph{e-mail}).
\end{enumerate}

\subsection{Para orientaçãopresencial}

\begin{enumerate}
  \item  A pessoa que solicita a informação entra em contato com a subsede de  referência para avaliar o melhor horário para o atendimento;
  \item  a informação é prestada na subsede.
\end{enumerate}

\section{Quanto custa?}

Gratuito.

\section{Quanto tempo leva?}

Não há prazo definido. Os atendimentos telefônicos e presenciais sãorealizados imediatamente pela equipe técnica, salvo quando esta seencontrar em atividade externa. Nesse caso, é oferecida à/ao solicitantea possibilidade de telefonar para outro território e pedir para a/otécnica/o da subsede atender a demanda.

Sendo de interesse falar diretamente com a/o técnica/o de referência,solicita-se o envio de \emph{e-mail} e/ou contato telefônico posterior.Para respostas por \emph{e-mail}, a COF leva em média de 1 a 7 dias pararetorno, a depender da complexidade da orientação solicitada.

\section{Como ter acesso ao atual \emph{status} do serviço?}

Não se aplica.

\section{Legislação relacionada}

\begin{itemize}
  \item  \href{http://www.planalto.gov.br/ccivil_03/leis/l5766.htm}{Lei Federal  nº~5.766/1971}, que cria o Conselho Federal e os Conselhos Regionais  de Psicologia;
  \item  \href{https://atosoficiais.com.br/cfp/resolucao-do-exercicio-profissional-n-10-2005-aprova-o-codigo-de-etica-profissional-do-psicologo}{Resolução  CFP nº~10/2005}, que institui o  \href{https://www.crpsp.org/uploads/pagina/289379/2j9LIMPLJ9jFr5YrK57HmAiBWjVMxdbe.pdf}{Código  de Ética Profissional da/o Psicóloga/o};
  \item  \href{https://transparencia.cfp.org.br/wp-content/uploads/sites/10/2017/07/Resolu\%C3\%A7\%C3\%A3o-CFP-010-2017-Politica-de-Orienta\%C3\%A7\%C3\%A3o-e-Fiscaliza\%C3\%A7\%C3\%A3o-dos-Conselhos.pdf}{Resolução  CFP nº~10/2017}, que institui a Política de Orientação e Fiscalização  do Sistema Conselhos de Psicologia
\end{itemize}

\chapter{\texorpdfstring{\emph{Consulta de legislaçãoprofissional}}{Consulta de legislação profissional}}


\section{O que é?}

Compilado de normativas (Resoluções do Conselho Federal de Psicologia edo Conselho Regional de Psicologia de São Paulo) que regulamentam aatuação profissional e os aspectos éticos e técnicos a serem observadosna prestação de serviços psicológicos em diferentes contextos.

\section{Quem pode utilizar esse serviço?}

O acesso é disponibilizado à categoria e à sociedade, não serestringindo apenas às/aos cidadãs/ãos brasileiras/os.

\section{Como solicitar o serviço?}

É possível ter acesso às normativas pelo\href{https://www.crpsp.org/legislacao/index}{\emph{site} do CRP~SP} oupelos \emph{site}s de atos oficiais do\href{https://atosoficiais.com.br/cfp}{CFP} e do\href{https://atosoficiais.com.br/crp06}{CRP}.

\section{Que informações ou documentos são necessários?}

Não é necessário enviar documentos nos casos de busca ativa nos\emph{site}s indicados.

\section{Quais são as etapas para a realização desse serviço?}

Acesso direto pelo\href{https://www.crpsp.org/legislacao/index}{\emph{site} do CRP~SP}, noqual as informações ficam organizadas por ano e são classificadas portipo de normativa (leis, resoluções, portarias, entre outras), ou pelos\emph{site}s de atos oficiais do\href{https://atosoficiais.com.br/cfp}{CFP} e do\href{https://atosoficiais.com.br/crp06}{CRP}.

\section{Quanto custa?}

Gratuito.

\section{Quanto tempo leva?}

Acesso imediato.

\section{Como ter acesso ao atual \emph{status} do serviço?}

Não se aplica.

\section{Legislação relacionada}

\begin{itemize}
  \item  \href{http://www.planalto.gov.br/ccivil_03/leis/l5766.htm}{Lei Federal  nº~5.766/1971}, que cria o Conselho Federal e os Conselhos Regionais  de Psicologia;
  \item  \href{https://transparencia.cfp.org.br/wp-content/uploads/sites/10/2017/07/Resolu\%C3\%A7\%C3\%A3o-CFP-010-2017-Politica-de-Orienta\%C3\%A7\%C3\%A3o-e-Fiscaliza\%C3\%A7\%C3\%A3o-dos-Conselhos.pdf}{Resolução  CFP nº~10/2017}, que institui a Política de Orientação e Fiscalização  do Sistema Conselhos de Psicologia
\end{itemize}

\section{Serviços correlatos}

Não se aplica.

\chapter{\texorpdfstring{\emph{Apresentação dequeixa}}{Apresentação de queixa}}


\section{O que é?}

É possível apresentar uma queixa anônima por \emph{e-mail} oupessoalmente em todas as subsedes do CRP~SP (confira os endereços\href{https://www.crpsp.org/sede/index}{no \emph{site}}) para análise eprovidências cabíveis.

\section{Quem pode utilizar esse serviço?}

O acesso é disponibilizado a qualquer cidadã/ão ou instituição.

\section{Como solicitar o serviço?}

A queixa anônima deve ser enviada preferencialmente para o \emph{e-mail}\href{mailto:orientacao@crpsp.org.br}{\nolinkurl{orientacao@crpsp.org.br}}.A pessoa que realizar uma queixa anônima ou sob anonimato não seráconsiderada parte e não será informada dos encaminhamentos realizados,uma vez que os procedimentos tramitam em sigilo. O anonimato serágarantido, excetuando-se solicitação judicial e/ou casos previstos emlei.

\section{Que informações ou documentos são necessários?}

É importante que a pessoa apresente dados detalhados que indiquem opossível dano causado pela/o profissional e/ou instituição, com aidentificação das/os psicólogas/os envolvidas/os.

\section{Quais são as etapas para a realização desse serviço?}

\begin{enumerate}
  \item  A Comissão de Orientação e Fiscalização verifica se há indícios de  irregularidade na conduta ética profissional apresentada como queixa;
\end{enumerate}

\begin{enumerate}
  \item  após análise, podem ser adotadas as seguintes ações:

        \begin{enumerate}
          \item    realização de atividades de fiscalização;
          \item    oferta de orientação à/ao psicóloga/o;
          \item    celebração de termo de ajustamento de conduta (TAC);
          \item    apresentação de requerimento de ofício (tendo a COF como    representante);
          \item    arquivamento da queixa, caso não se constatem indícios de    irregularidade do ponto de vista ético;
          \item    outras que julgar pertinentes, considerando as atribuições legais da    autarquia.
        \end{enumerate}
\end{enumerate}

\section{Quanto custa?}

Gratuito.

\section{Quanto tempo leva?}

Não há prazo normatizado em legislação.

\section{Como ter acesso ao atual \emph{status} do serviço?}

No caso de queixas anônimas, os procedimentos tramitam sob sigilo, sendopossível apenas o acompanhamento das representações.

\section{Legislação relacionada}

\begin{itemize}
  \item  \href{http://www.planalto.gov.br/ccivil_03/leis/l5766.htm}{Lei Federal  nº~5.766/1971}, que cria o Conselho Federal e os Conselhos Regionais  de Psicologia;
  \item  \href{https://transparencia.cfp.org.br/wp-content/uploads/sites/10/2017/07/Resolu\%C3\%A7\%C3\%A3o-CFP-010-2017-Politica-de-Orienta\%C3\%A7\%C3\%A3o-e-Fiscaliza\%C3\%A7\%C3\%A3o-dos-Conselhos.pdf}{Resolução  CFP nº~10/2017}, que institui a Política de Orientação e Fiscalização  do Sistema Conselhos de Psicologia
\end{itemize}

\chapter{\texorpdfstring{\emph{Representação}}{Representação}}


\section{O que é?}

Qualquer pessoa que se sinta lesada pode fazer uma representação paraque seja apurada a conduta ética da/o psicóloga/o ou pessoa jurídicainscritas/os no CRP~SP.

Apresentada a representação, o CRP~SP poderá decidir por arquivamento,envio para mediação ou instauração de processo investigativo oudisciplinar em face da/o psicóloga/o.

\section{Quem pode utilizar esse serviço?}

Qualquer pessoa interessada ou que tenha sido prejudicada pelo serviçoda/o psicóloga/o.

\section{Como solicitar o serviço?}

A pessoa interessada deve preencher o\hyperref[]{formulário derepresentação} e enviá-lo exclusivamente para o \emph{e-mail}\href{mailto:representacao@crpsp.org.br}{\nolinkurl{representacao@crpsp.org.br}},junto com um documento escrito e assinado em que se identifique edescreva os fatos conhecidos por meio das diversas mídias, peloresultado de visitas de fiscalização, por declarações e manifestaçõespúblicas ou de outras formas. Também podem ser anexados outrosdocumentos que sustentem a alegação.

Mais informações estão disponíveis no\href{https://www.crpsp.org/pagina/view/50}{\emph{site} do CRP~SP} e noLivro II, Título I, Capítulo I do\href{https://atosoficiais.com.br/cfp/resolucao-de-fiscalizacao-e-orientacao-n-11-2019-institui-o-codigo-de-processamento-disciplinar?origin=instituicao&q=011/2019}{Códigode Processamento Disciplinar}.

\section{Que informações ou documentos são necessários?}

\begin{itemize}
  \item  \hyperref[]{Formulário de  representação} preenchido e assinado;
  \item  documento de identificação (RG, CNH etc. --- necessário apenas para  pessoas que não possuírem \hl{assinatura digital certificada});
  \item  documento escrito e assinado em que se identifique e descreva os fatos  conhecidos por meio das diversas mídias, pelo resultado de visitas de  fiscalização, por declarações e manifestações públicas ou de outras  formas;
  \item  outros documentos que sustentem a alegação
\end{itemize}

\section{Quais são as etapas para a realização desse serviço?}

\begin{enumerate}
  \item  A pessoa que fará a representação elabora documento em que se  identifica e narra os fatos que a levaram a acionar o CRP~SP, preenche  e assina o \hyperref[]{formulário  de representação};
\end{enumerate}

\begin{enumerate}
  \item  formulário e documento narrando os fatos são enviados para o  \emph{e-mail}  \href{mailto:representacao@crpsp.org.br}{\nolinkurl{representacao@crpsp.org.br}},  endereçados à Presidência do Conselho;
  \item  \hl{a Comissão de Ética delibera sobre a representação e toma as  providências necessárias}.
\end{enumerate}

Embora seja um procedimento administrativo extrajudicial, arepresentação no CRP~SP possui um trâmite análogo ao processo judicial,que envolve várias etapas. Em todas as fases processuais, é possívelocorrer a mediação.

\section{Quanto custa?}

Gratuito.

\section{Quanto tempo leva?}

Se a representação culminar na instauração de um processo investigativoe disciplinar, o trâmite todo pode durar até cinco anos.

\section{Como ter acesso ao atual \emph{status} do serviço?}

Após apresentada a representação, o usuário será informado de cada etapaprocessual. Também é possível consultar o andamento ou obter cópiadigital do processo por meio do \emph{e-mail}\href{mailto:etica04@crpsp.org.br}{\nolinkurl{etica04@crpsp.org.br}}

\section{Legislação relacionada}

\begin{itemize}
  \item  \href{http://www.planalto.gov.br/ccivil_03/leis/l5766.htm}{Lei Federal  nº~5.766/1971}, que cria o Conselho Federal e os Conselhos Regionais  de Psicologia;
  \item  \href{https://atosoficiais.com.br/cfp/resolucao-do-exercicio-profissional-n-10-2005-aprova-o-codigo-de-etica-profissional-do-psicologo}{Resolução  CFP nº~10/2005}, que institui o  \href{https://www.crpsp.org/uploads/pagina/289379/2j9LIMPLJ9jFr5YrK57HmAiBWjVMxdbe.pdf}{Código  de Ética Profissional da/o Psicóloga/o};
  \item  \href{https://atosoficiais.com.br/lei/codigo-de-processamento-disciplinar-cfp}{Resolução  CFP nº~11/2019}, que institui o Código de Processamento Disciplinar
\end{itemize}

\section{Serviços correlatos}

\begin{itemize}
  \item  \hyperref[]{Apresentação de queixa};
  \item  \hyperref[]{orientações  sobre o Código de Ética e legislação profissional};
  \item  \hyperref[]{consulta de  legislação profissional}
\end{itemize}

\chapter{\texorpdfstring{\emph{Orientação/fiscalização de pessoasjurídicasrecém-inscritas}}{Orientação/fiscalização de pessoas jurídicas recém-inscritas}}


\section{O que é?}

A Comissão de Orientação e Fiscalização (COF) do CRP~SP orienta pessoasjurídicas (PJ) recém-inscritas por meio de ofícios orientativos e/oudurante ações rotineiras de fiscalização.

\section{Quem pode utilizar esse serviço?}

Ao concluir a inscrição no CRP~SP, todas as recebem ofício de orientaçãodestinado às/aos Responsáveis Técnicas/os do serviço.

Já os fluxos de fiscalização são estabelecidos a partir da organizaçãointerna da Comissão de Orientação e Fiscalização.

\section{Como solicitar o serviço?}

Não é necessário solicitar serviços. Tanto as ações de fiscalizaçãoquanto o envio de ofícios de orientação são procedimentos obrigatóriospara todas as empresas inscritas, conforme a\href{https://transparencia.cfp.org.br/wp-content/uploads/sites/10/2017/07/Resolu\%C3\%A7\%C3\%A3o-CFP-010-2017-Politica-de-Orienta\%C3\%A7\%C3\%A3o-e-Fiscaliza\%C3\%A7\%C3\%A3o-dos-Conselhos.pdf}{ResoluçãoCFP nº~10/2017}.

\section{Que informações ou documentos são necessários?}

\begin{itemize}
  \item  \hyperref[]{Inscrição  profissional de pessoa jurídica}
\end{itemize}

\section{Quais são as etapas para a realização desse serviço?}

\begin{enumerate}
  \item  A inscrição profissional de pessoa jurídica é deferida;
\end{enumerate}

\begin{enumerate}
  \item  a COF envia ofícios orientativos e/ou procede à fiscalização de  rotina.
\end{enumerate}

\section{Quanto custa?}

Gratuito.

\section{Quanto tempo leva?}

As orientações precisam ser enviadas em até 120 dias da conclusão dopedido de inscrição da empresa. No CRP~SP, o procedimento é finalizadojunto com a confirmação da inscrição.

As ações de fiscalização são definidas pela COF, após análise do objetosocial da empresa. Não existe prazo determinado para a execução dessasações.

\section{Como ter acesso ao atual \emph{status} do serviço?}

Não se aplica.

\section{Legislação relacionada}

\begin{itemize}
  \item  \href{http://www.planalto.gov.br/ccivil_03/leis/l5766.htm}{Lei Federal  nº~5.766/1971}, que cria o Conselho Federal e os Conselhos Regionais  de Psicologia;
  \item  \href{https://atosoficiais.com.br/cfp/resolucao-do-exercicio-profissional-n-10-2005-aprova-o-codigo-de-etica-profissional-do-psicologo}{Resolução  CFP nº~10/2005}, que institui o  \href{https://www.crpsp.org/uploads/pagina/289379/2j9LIMPLJ9jFr5YrK57HmAiBWjVMxdbe.pdf}{Código  de Ética Profissional da/o Psicóloga/o};
  \item  \href{https://transparencia.cfp.org.br/wp-content/uploads/sites/10/2017/07/Resolu\%C3\%A7\%C3\%A3o-CFP-010-2017-Politica-de-Orienta\%C3\%A7\%C3\%A3o-e-Fiscaliza\%C3\%A7\%C3\%A3o-dos-Conselhos.pdf}{Resolução  CFP nº~10/2017}, que institui a Política de Orientação e Fiscalização  do Sistema Conselhos de Psicologia.
  \item  \href{https://atosoficiais.com.br/cfp/resolucao-administrativa-financeira-n-16-2019-altera-o-manual-de-procedimentos-administrativo-e-financeiro-do-sistema-conselhos-de-psicologia-anexo-da-resolucao-cfp-n-202018-a-resolucao-cfp-n-03-2007-a-resolucao-cfp-n-16-2019-e-da-outras-providencias}{Resolução  CFP nº~16/2019}, que normatiza o registro e o cadastro de pessoas  jurídicas;
  \item  \href{https://atosoficiais.com.br/cfp/resolucao-administrativa-financeira-n-8-2023-altera-o-manual-de-procedimentos-administrativo-e-financeiro-do-sistema-conselhos-de-psicologia-anexo-da-resolucao-cfp-n-202018-a-resolucao-cfp-n-03-2007-a-resolucao-cfp-n-16-2019-e-da-outras-providencias?origin=instituicao}{Resolução  CFP nº~08/2023}, que altera o  \href{https://transparencia.cfp.org.br/wp-content/uploads/sites/7/2020/05/Manual-de-Procedimentos-Administrativos-e-Financeiro-CFP-RES-20.2018.pdf}{\emph{Manual  de procedimentos administrativos e financeiros}}
\end{itemize}

\section{Serviços correlatos}

\begin{itemize}
  \item  \hyperref[]{Inscrição  de pessoa jurídica no CRP~SP}
\end{itemize}

\section{Anexos}

\chapter{\texorpdfstring{\emph{Requerimento de cancelamento deinscrição (pessoafísica)}}{Requerimento de cancelamento de inscrição (pessoa física)}}

\href{file:///C:/Users/angelo.cuissi/Documents/cuissi/caderno/revisão/202406/carta-de-servicos/carta-preparacao-de-original/novos_requerimentos/PF-cancelamento-inscricao_p.pdf}{novos\_requerimentos\textbackslash PF-cancelamento-inscricao\_p.pdf}

\chapter{\texorpdfstring{\emph{Requerimento de registro deespecialista}}{Requerimento de registro de especialista}}

\href{file:///C:/Users/angelo.cuissi/Documents/cuissi/caderno/revisão/202406/carta-de-servicos/carta-preparacao-de-original/PF-requerimento-especialista.pdf}{Requerimento}

\chapter{\texorpdfstring{\emph{Requerimento de inscrição (pessoajurídica)}}{Requerimento de inscrição (pessoa jurídica)}}

\href{file:///C:/Users/angelo.cuissi/Documents/cuissi/caderno/revisão/202406/carta-de-servicos/carta-preparacao-de-original/PJ-inscricao.pdf}{PJ-inscricao.pdf}

\chapter{\texorpdfstring{\emph{Termo de responsabilidade técnica (RT)e autorização para envio de comunicaçõeseletrônicas}}{Termo de responsabilidade técnica (RT) e autorização para envio de comunicações eletrônicas}}

\href{file:///C:/Users/angelo.cuissi/Documents/cuissi/caderno/revisão/202406/carta-de-servicos/carta-preparacao-de-original/PJ-termo-RT.pdf}{PJ-termo-RT.pdf}

\chapter{\texorpdfstring{\emph{Declaração institucional para garantiado amplo e livre exercício profissional por parte dos profissionais dePsicologia e Responsáveis Técnicos de Psicologia de acordo com o Códigode Ética Profissional da/o Psicóloga/o e legislaçãocabível}}{Declaração institucional para garantia do amplo e livre exercício profissional por parte dos profissionais de Psicologia e Responsáveis Técnicos de Psicologia de acordo com o Código de Ética Profissional da/o Psicóloga/o e legislação cabível}}

\href{file:///C:/Users/angelo.cuissi/Documents/cuissi/caderno/revisão/202406/carta-de-servicos/carta-preparacao-de-original/PJ-garantia-autorizacao.pdf}{PJ-garantia-autorizacao.pdf}

\chapter{\texorpdfstring{\emph{Solicitação de cancelamento/alteraçãode registro de pessoajurídica}}{Solicitação de cancelamento/alteração de registro de pessoa jurídica}}

\href{file:///C:/Users/angelo.cuissi/Documents/cuissi/caderno/revisão/202406/carta-de-servicos/carta-preparacao-de-original/PJ-cancelamento-alteracao.pdf}{PJ-cancelamento-alteracao.pdf}

\chapter{\texorpdfstring{\emph{Formulário de declaração dedesligamento de ResponsávelTécnica/o}}{Formulário de declaração de desligamento de Responsável Técnica/o}}

\href{file:///C:/Users/angelo.cuissi/Documents/cuissi/caderno/revisão/202406/carta-de-servicos/carta-preparacao-de-original/PJ-desligamento-RT.pdf}{PJ-desligamento-RT.pdf}

\chapter{\texorpdfstring{\emph{Formulário de informações sobre asatividades da pessoa jurídica para análise do pedido decancelamento}}{Formulário de informações sobre as atividades da pessoa jurídica para análise do pedido de cancelamento}}

\href{file:///C:/Users/angelo.cuissi/Documents/cuissi/caderno/revisão/202406/carta-de-servicos/carta-preparacao-de-original/PJ-informacoes-atividades.pdf}{PJ-informacoes-atividades.pdf}

\chapter{\texorpdfstring{\emph{Solicitação de renovação decertificado}}{Solicitação de renovação de certificado}}

\href{file:///C:/Users/angelo.cuissi/Documents/cuissi/caderno/revisão/202406/carta-de-servicos/carta-preparacao-de-original/PJ-renovacao.pdf}{PJ-renovacao.pdf}

\chapter{\texorpdfstring{\emph{Solicitação de isenção dos débitos depessoa jurídica por alteraçãocontratual}}{Solicitação de isenção dos débitos de pessoa jurídica por alteração contratual}}

\href{file:///C:/Users/angelo.cuissi/Documents/cuissi/caderno/revisão/202406/carta-de-servicos/carta-preparacao-de-original/PJ-isencao.pdf}{Formulário}

\chapter{\texorpdfstring{\emph{Formulário derepresentação}}{Formulário de representação}}

\href{file:///C:/Users/angelo.cuissi/Documents/cuissi/caderno/revisão/202406/carta-de-servicos/carta-preparacao-de-original/formulario-de-representacao.pdf}{formulario-de-representacao.pdf}

\end{document}